% Source: Gerzensee "Probability Theory" slides 3-32 (probability space and the
% four properties of P; conditional probability; Bayes Theorems 1 and 2 and the
% partition form; the medical-test, airport-carousel, market-microstructure and
% political-poll examples with their numerical answers; random variables, pmf and
% expectation in the two equivalent forms; conditional expectation given an event
% and given a random variable; Theorems 2.4 and 2.5, the two laws of iterated
% expectations; conditional expectation as best predictor with the full proof;
% normal-normal updating in variance form and precision form, the completing-the-
% square derivation, and the multiple-observation version; the closing list of
% applications).
% Supplementary: Fukuda (2026), Sec. 2.1-2.5 (pp. 57-82; Def. 2.1-2.13,
% Prop. 2.1, 2.7, Remark 2.1-2.5, Thm 2.4-2.5, Exercises 2.1, 2.11, 2.14, 2.15,
% 2.17-2.20) and Sec. 9.6 (pp. 323-332; best predictor, best linear predictor,
% the bivariate normal case, Thm 9.5-9.6, Markov and Chebyshev).
% External references actually cited in this chapter: Brockwell-Davis (2016),
% Gelman et al. (2013), Grossman-Stiglitz (1980), Kyle (1985), Le Gall (2022),
% O'Hara (1998).

\chapter[Probability and Conditional Expectation]{Probability, Conditional Expectation and Bayesian Updating}
\label{ch:probability}

\section{Probability Spaces}\label{sec:prob-spaces}

Uncertainty enters an economic model as a set of states, a rule assigning
likelihoods to statements about the state, and functions of the state that the
agent cares about. The three objects are the sample space, the probability
measure and the random variables, and the first difficulty is that the second
cannot in general be defined on all statements.

\begin{definition}[$\sigma$-algebra]\label{def:sigma-algebra}
	Let $\Omega$ be a set. A collection $\mathcal{F}$ of subsets of $\Omega$ is a
	\dfnas{sigma-algebra@$\sigma$-algebra}{$\sigma$-algebra} on $\Omega$ if
	\begin{enumerate}[(i)]
		\item $\Omega\in\mathcal{F}$;
		\item $E\in\mathcal{F}$ implies $E^{c}\in\mathcal{F}$;
		\item $E_n\in\mathcal{F}$ for every $n\in\N$ implies $\bigcap_{n\in\N}E_n\in\mathcal{F}$.
	\end{enumerate}
	The pair $(\Omega,\mathcal{F})$ is then a \dfnas{measurable space}{measurable space} and
	the members of $\mathcal{F}$ are \dfnas{event}{events}.
	\nidx{sigmaalgebra}{$\mathcal{F}$, $\sigma$-algebra of events}
\end{definition}

\begin{proposition}[Closure properties]\label{prop:sigma-algebra-basics}
	Let $\mathcal{F}$ be a $\sigma$-algebra on $\Omega$. Then $\emptyset\in\mathcal{F}$ and
	$\bigcup_{n\in\N}E_n\in\mathcal{F}$ whenever every $E_n\in\mathcal{F}$; consequently $\mathcal{F}$ is
	closed under finite unions, finite intersections and set differences.
\end{proposition}

\begin{proof}
	$\emptyset=\Omega^{c}\in\mathcal{F}$ by (i) and (ii). For unions, De Morgan's law
	(\autoref{prop:demorgan}) gives
	$\bigcup_nE_n=\bigl(\bigcap_nE_n^{c}\bigr)^{c}$, and each $E_n^{c}\in\mathcal{F}$ by
	(ii), so the intersection lies in $\mathcal{F}$ by (iii) and its complement by (ii).
	Finite operations follow by padding the sequence with $\Omega$ or $\emptyset$,
	and $E\setminus F=E\cap F^{c}$.
\end{proof}

An event is a statement about the experiment admitting a yes-or-no answer, and
the three axioms are the closure of such statements under negation and countable
conjunction. The restriction to \emph{countable} conjunction is what makes the
definition non-trivial: closure under arbitrary unions would force $\mathcal{F}$
to contain every subset of $\Omega$, since every set is the union of its
singletons.

That the whole power set is inadmissible is a statement about a
\emph{particular} measure, not about probability measures in general. The Dirac
measure $\delta_{\omega_0}(E)=\one\{\omega_0\in E\}$ is a countably additive
probability measure on $\pow(\Omega)$ for any $\Omega$ whatsoever. What fails is
the following: there is no countably additive probability measure on
$\pow\bigl([0,1]\bigr)$ that is invariant under translation modulo $1$ and
assigns to each interval its length. The standard witness is a \dfnas{Vitali set}{Vitali
	set} --- a set of representatives of the cosets of $\Q$ in
$\R$, whose existence uses the axiom of choice --- for which countable additivity
and translation invariance force the measure of $[0,1]$ to be either $0$ or
$\infty$. Lebesgue measure therefore has to be restricted to the Borel or
Lebesgue $\sigma$-algebra, and every uniform-distribution argument in economics
inherits that restriction. See \textcite[Ch.~1]{legall2022measure}.
\autoref{ch:topology}
recorded the formal parallel with \autoref{def:topological-space}: a topology is
closed under arbitrary unions and finite intersections, a $\sigma$-algebra under
countable operations of both kinds and under complementation, and neither
contains the other.

\begin{definition}[Probability measure and probability space]\label{def:probability-measure}
	Let $(\Omega,\mathcal{F})$ be a measurable space. A function $\Prob\colon\mathcal{F}\to[0,1]$ is a
	\dfnas{probability measure}{probability measure} if
	\begin{enumerate}[(i)]
		\item $\Prob(\Omega)=1$;
		\item ($\sigma$-additivity) $\Prob\bigl(\bigcup_{n\in\N}E_n\bigr)
			      =\sum_{n=1}^{\infty}\Prob(E_n)$ whenever $(E_n)_{n\in\N}$ is pairwise
		      disjoint.
	\end{enumerate}
	The triple $(\Omega,\mathcal{F},\Prob)$ is a \dfnas{probability space}{probability space}.
	\nidx{probmeasure}{$\Prob$, probability measure}
\end{definition}

\begin{proposition}[Elementary properties]\label{prop:probability-basics}
	Let $(\Omega,\mathcal{F},\Prob)$ be a probability space and $E,F\in\mathcal{F}$.
	\begin{enumerate}[(i)]
		\item $\Prob(\emptyset)=0$.
		\item (Monotonicity) $E\subseteq F$ implies $\Prob(E)\leq\Prob(F)$.
		\item (Finite additivity) $E\cap F=\emptyset$ implies
		      $\Prob(E\cup F)=\Prob(E)+\Prob(F)$; in particular
		      $\Prob(E)+\Prob(E^{c})=1$.
		\item (Inclusion--exclusion)
		      $\Prob(E\cup F)=\Prob(E)+\Prob(F)-\Prob(E\cap F)$, hence
		      $\Prob(E\cup F)\leq\Prob(E)+\Prob(F)$.
		\item (Continuity from below) If $E_1\subseteq E_2\subseteq\cdots$ and
		      $E=\bigcup_nE_n$, then $\Prob(E_n)\to\Prob(E)$.
		\item (Continuity from above) If $F_1\supseteq F_2\supseteq\cdots$ and
		      $F=\bigcap_nF_n$, then $\Prob(F_n)\to\Prob(F)$.
	\end{enumerate}
\end{proposition}

\begin{proof}
	(i) Suppose $\Prob(\emptyset)=a>0$ and take $E_n=\emptyset$ for every $n$. These
	are pairwise disjoint, so $\sigma$-additivity gives
	$a=\Prob(\emptyset)=\sum_{n=1}^{\infty}a=\infty$, contradicting
	$\Prob\leq1$.

	(ii) Put $E_1=E$, $E_2=F\setminus E$ and $E_n=\emptyset$ for $n\geq3$; these are
	disjoint with union $F$, so
	$\Prob(F)=\Prob(E)+\Prob(F\setminus E)\geq\Prob(E)$ by (i) and non-negativity.

	(iii) Put $E_1=E$, $E_2=F$, $E_n=\emptyset$ for $n\geq3$. Taking $F=E^{c}$ gives
	$\Prob(E)+\Prob(E^{c})=\Prob(\Omega)=1$.

	(iv) Decompose $E\cup F$ into the three disjoint events
	$E\setminus F$, $E\cap F$, $F\setminus E$ and apply (iii) three times:
	$\Prob(E\cup F)=\Prob(E\setminus F)+\Prob(E\cap F)+\Prob(F\setminus E)$, while
	$\Prob(E)=\Prob(E\setminus F)+\Prob(E\cap F)$ and similarly for $F$; adding the
	last two and subtracting $\Prob(E\cap F)$ gives the identity. The inequality
	follows since $\Prob(E\cap F)\geq0$.

	(v) Put $A_1\eqdef E_1$ and $A_n\eqdef E_n\setminus E_{n-1}$ for $n\geq2$. The
	$A_n$ are pairwise disjoint, $\bigcup_{k\leq n}A_k=E_n$ and
	$\bigcup_kA_k=E$, so $\sigma$-additivity gives
	$\Prob(E)=\sum_{k\geq1}\Prob(A_k)=\lim_n\sum_{k\leq n}\Prob(A_k)=\lim_n\Prob(E_n)$.

	(vi) Apply (v) to $E_n\eqdef\comp{F_n}$, which increase to $\comp{F}$, and use
	(iii).
\end{proof}

\begin{remark}[The countable case]\label{rem:discrete-probability}
	When $\Omega$ is at most countable one may take $\mathcal{F}=\pow(\Omega)$ and specify
	$\Prob$ by the numbers $\Prob(\{\omega\})$, $\omega\in\Omega$, which are
	non-negative and sum to $1$; every event then has
	$\Prob(E)=\sum_{\omega\in E}\Prob(\{\omega\})$ and $\sigma$-additivity is
	automatic. The four properties listed on
	\autocite[slide 4]{fukuda2026probability} --- $\Prob(\emptyset)=0$,
	$\Prob(\Omega)=1$, monotonicity and additivity on disjoint events --- are exactly
	\autoref{def:probability-measure} together with
	\autoref{prop:probability-basics} specialised to this case. The material of
	\autoref{sec:prob-bayes} and \autoref{sec:prob-conditional} is developed in this
	countable setting, where every expectation is a sum and conditioning is
	carried out only on events of strictly positive probability. A countable
	space still admits null events --- a point mass $\Prob(\{\omega\})=0$ is
	allowed by the definition, as is a countable union of such points --- so the
	restriction is a choice made in those sections and not a property of
	countability. The passage to general probability spaces is not a
	matter of rewriting sums as integrals: conditioning on an event of probability
	zero has no meaning under \autoref{def:conditional-probability}, and the object
	that replaces it is defined through a $\sigma$-algebra and is unique only up to
	a null set. \autoref{subsec:prob-general-conditioning} carries this out, and it
	is what licenses the continuous-state calculations of
	\autoref{sec:prob-normal} and of \autoref{ch:dynamic}.
\end{remark}

\section{Random Variables and Expectation}\label{sec:prob-rv}

\begin{definition}[Random variable]\label{def:random-variable}
	Let $(\Omega,\mathcal{F},\Prob)$ be a probability space. A \dfnas{random variable}{random
		variable} is a function $X\colon\Omega\to\R$ that is
	\dfnas{measurable function}{measurable}: $X^{-1}(B)\in\mathcal{F}$ for every Borel set
	$B\subseteq\R$. When $\Omega$ is at most countable and $\mathcal{F}=\pow(\Omega)$ the
	measurability requirement is vacuous. Writing $X(\Omega)=\{x_i\}_{i}$ for the
	set of realisations, the \dfnas{probability mass
		function}{probability mass function} of $X$ is
	\begin{equation}\label{eq:pmf}
		f(x)\eqdef\Prob\bigl(\{\omega\in\Omega : X(\omega)=x\}\bigr)
		=\Prob\bigl(X^{-1}(\{x\})\bigr),\qquad x\in X(\Omega),
	\end{equation}
	and one writes $p_i\eqdef f(x_i)$ and $[X=x]\eqdef X^{-1}(\{x\})$.
	\nidx{randomvariable}{$X$, $Y$, random variables}
\end{definition}

Measurability is the requirement that every statement about the value of $X$ be
an event, so that its probability is defined. It is the probabilistic
counterpart of continuity: \autoref{thm:continuity-preimage} characterised
continuity by $f^{-1}(\text{open})$ open, and
\autoref{rem:continuity-measurability} recorded that measurability replaces
``open'' by ``Borel'' and ``open'' by ``in $\mathcal{F}$''. On a countable sample space the
condition is empty, which is why the slides define a random variable simply as a
function \autocite[slide 15]{fukuda2026probability}.

\begin{definition}[Expectation]\label{def:expectation}
	For a random variable $X$ on an at most countable probability space with
	$X(\Omega)=\{x_i\}_{i=1}^{n}$, the \dfnas{expectation}{expectation} is
	\begin{equation}\label{eq:expectation}
		\E[X]\eqdef\sum_{i=1}^{n}p_ix_i
		=\sum_{\omega\in\Omega}\Prob(\{\omega\})X(\omega),
	\end{equation}
	the two expressions being equal because the events $[X=x_i]$ partition $\Omega$.
	\nidx{expectation}{$\E[X]$, expectation}
\end{definition}

\begin{proof}[Proof of the second equality in \eqref{eq:expectation}]
	Group the sum over $\Omega$ by the value of $X$:
	\[
		\sum_{\omega\in\Omega}\Prob(\{\omega\})X(\omega)
		=\sum_{i=1}^{n}\ \sum_{\omega\,:\,X(\omega)=x_i}\Prob(\{\omega\})X(\omega)
		=\sum_{i=1}^{n}x_i\sum_{\omega\,:\,X(\omega)=x_i}\Prob(\{\omega\})
		=\sum_{i=1}^{n}x_ip_i,
	\]
	using \eqref{eq:pmf} in the last step. Rearrangement is legitimate because the
	family $\bigl([X=x_i]\bigr)_{i}$ partitions $\Omega$ and, for a countable
	$\Omega$, absolute convergence of $\sum_{\omega}\Prob(\{\omega\})\abs{X(\omega)}$
	is exactly the integrability of $X$.
\end{proof}

\begin{proposition}[Properties of expectation]\label{prop:expectation-properties}
	Let $X,Y$ be integrable random variables and $a,b\in\R$.
	\begin{enumerate}[(i)]
		\item (Linearity) $\E[aX+bY]=a\E[X]+b\E[Y]$.
		\item (Monotonicity) $X\leq Y$ pointwise implies $\E[X]\leq\E[Y]$; and
		      $X\geq0$ with $\E[X]=0$ implies $\Prob(X=0)=1$.
		\item (Change of variables) $\E[g(X)]=\sum_{x\in X(\Omega)}g(x)f(x)$ for any
		      $g\colon\R\to\R$ for which the sum converges absolutely.
		\item (Independence) If $X$ and $Y$ are independent, meaning
		      $\Prob([X\in A]\cap[Y\in B])=\Prob(X\in A)\Prob(Y\in B)$ for all Borel
		      $A,B$, then $\E[XY]=\E[X]\E[Y]$.
	\end{enumerate}
\end{proposition}

\begin{proof}
	(i)--(iii) are rearrangements of absolutely convergent sums, using
	\eqref{eq:expectation} in its second form for (i) and its first form for (iii).
	For the second half of (ii), if $\Prob(X>0)>0$ then $\Prob(X>1/k)>0$ for some
	$k$ by countable additivity applied to $[X>0]=\bigcup_k[X>1/k]$, and then
	$\E[X]\geq(1/k)\Prob(X>1/k)>0$. For (iv), independence makes the joint mass
	function factor, $f_{X,Y}(x,y)=f_X(x)f_Y(y)$, and
	$\E[XY]=\sum_{x,y}xyf_X(x)f_Y(y)=\bigl(\sum_xxf_X(x)\bigr)\bigl(\sum_yyf_Y(y)\bigr)$.
\end{proof}

\begin{definition}[Variance and covariance]\label{def:variance}
	For square-integrable $X,Y$,
	\[
		\Var(X)\eqdef\E\bigl[(X-\E[X])^{2}\bigr],
		\qquad
		\Cov(X,Y)\eqdef\E\bigl[(X-\E[X])(Y-\E[Y])\bigr],
	\]
	\[
		\Corr(X,Y)\eqdef\frac{\Cov(X,Y)}{\sqrt{\Var(X)\Var(Y)}}\,,
		\qquad\text{defined when }\Var(X)\Var(Y)>0 .
	\]
	\nidx{variance}{$\Var(X)$, $\Cov(X,Y)$, variance and covariance}
\end{definition}

\begin{proposition}[The variance identity]\label{prop:variance-identity}
	For square-integrable $X$ and $Y$,
	\begin{equation}\label{eq:variance-identity}
		\Var(X)=\E[X^{2}]-\bigl(\E[X]\bigr)^{2},
		\qquad
		\Cov(X,Y)=\E[XY]-\E[X]\E[Y] .
	\end{equation}
	More generally, for any $c\in\R$,
	$\E\bigl[(X-c)^{2}\bigr]=\Var(X)+\bigl(\E[X]-c\bigr)^{2}$, so
	$c\mapsto\E[(X-c)^{2}]$ is minimised uniquely at $c=\E[X]$.
\end{proposition}

\begin{proof}
	Write $\mu\eqdef\E[X]$ and expand, using
	\autoref{prop:expectation-properties}(i) and the fact that $\mu$ is a constant:
	\[
		\E\bigl[(X-c)^{2}\bigr]
		=\E\bigl[(X-\mu)^{2}\bigr]+2(\mu-c)\underbrace{\E[X-\mu]}_{=0}+(\mu-c)^{2}
		=\Var(X)+(\mu-c)^{2}.
	\]
	Taking $c=0$ gives the first identity in \eqref{eq:variance-identity}; the
	covariance identity is the same expansion applied to the product. The
	minimisation claim is immediate since $(\mu-c)^{2}\geq0$ with equality only at
	$c=\mu$.
\end{proof}

\begin{proposition}[Variance of a sum]\label{prop:variance-independent}
	Let $X_1,\dots,X_n$ be square-integrable and pairwise independent. Then
	\[
		\Var\Bigl(\sum_{i=1}^{n}X_i\Bigr)=\sum_{i=1}^{n}\Var(X_i),
		\qquad\text{and}\qquad
		\Var(aX+b)=a^{2}\Var(X)\ \text{ for }a,b\in\R .
	\]
	In particular, if the $X_i$ are identically distributed with variance
	$\sigma^{2}$, the sample mean $\bar X_n$ has $\Var(\bar X_n)=\sigma^{2}/n$.
\end{proposition}

\begin{proof}
	Expanding the square and using linearity,
	$\Var\bigl(\sum_iX_i\bigr)=\sum_i\Var(X_i)+\sum_{i\neq j}\Cov(X_i,X_j)$.
	Pairwise independence gives $\E[X_iX_j]=\E[X_i]\E[X_j]$ by
	\autoref{prop:expectation-properties}(iv), so
	$\Cov(X_i,X_j)=0$ for $i\neq j$ by \eqref{eq:variance-identity}. The scaling
	identity is $\E[(aX+b-a\E[X]-b)^{2}]=a^{2}\E[(X-\E[X])^{2}]$, and the last
	claim combines the two with $\bar X_n=n^{-1}\sum_iX_i$.
\end{proof}

The last sentence is the reason the mean, and not some other centre, is the
object that appears throughout asset pricing and econometrics: it is the constant
that best predicts $X$ in mean square. \autoref{thm:best-predictor} extends the
statement from constants to functions of a conditioning variable, and
\autoref{prop:best-linear-predictor} to affine functions.

\begin{theorem}[Jensen's inequality]\label{thm:jensen}
	Let $I\subseteq\R$ be an interval, $\varphi\colon I\to\R$ convex, and $X$ an
	integrable random variable with $\Prob(X\in I)=1$ and $\E[X]$ in the interior of
	$I$. Then $\E[\varphi(X)]\geq\varphi(\E[X])$, with equality if and only if
	$\varphi$ is affine on the smallest interval carrying all the mass of $X$. The
	inequality reverses for concave $\varphi$.
\end{theorem}

\begin{proof}
	Put $\mu\eqdef\E[X]$. By \autoref{prop:convex-epigraph} the set $\epi\varphi$ is
	convex, and $(\mu,\varphi(\mu))$ lies on its boundary, so
	\autoref{thm:supporting-hyperplane} supplies a supporting hyperplane: there is
	$(a,b)\neq0$ with $a x+b t\geq a\mu+b\varphi(\mu)$ for every
	$(x,t)\in\epi\varphi$. Since $t$ may be taken arbitrarily large, $b\geq0$; and
	$b=0$ would force $ax\geq a\mu$ for all $x\in I$ with $a\neq0$, impossible
	because $\mu$ is interior to $I$. Dividing by $b>0$ and setting
	$s\eqdef-a/b$ gives the supporting line
	\begin{equation}\label{eq:supporting-line}
		\varphi(x)\geq\varphi(\mu)+s(x-\mu)\qquad\text{for all }x\in I .
	\end{equation}
	Substituting $X$ and taking expectations, monotonicity and linearity
	(\autoref{prop:expectation-properties}) give
	$\E[\varphi(X)]\geq\varphi(\mu)+s\E[X-\mu]=\varphi(\mu)$. Equality forces
	$\E\bigl[\varphi(X)-\varphi(\mu)-s(X-\mu)\bigr]=0$ with a non-negative
	integrand, hence by \autoref{prop:expectation-properties}(ii) the integrand
	vanishes almost surely: $\varphi$ agrees with the affine function
	\eqref{eq:supporting-line} on the support of $X$.
\end{proof}

\begin{eg}[Risk aversion is Jensen]\label{eg:jensen-risk}
	If $u$ is concave, \autoref{thm:jensen} gives
	$\E[u(W)]\leq u(\E[W])$: the agent prefers the expected value of the gamble to
	the gamble. The certainty equivalent $c$ defined by $u(c)=\E[u(W)]$ therefore
	satisfies $c\leq\E[W]$, and the risk premium $\E[W]-c$ is non-negative. The
	second-order expansion of \autoref{thm:arrow-pratt} quantified this premium as
	$\tfrac12A(\E[W])\Var(W)$ for small risks; Jensen gives the sign for risks of
	any size, with no differentiability and no smallness. Equality holds only if $u$
	is affine on the support of $W$, which is the definition of risk neutrality
	there.
\end{eg}

\section{Conditional Probability and Bayes' Rule}\label{sec:prob-bayes}

\begin{definition}[Conditional probability]\label{def:conditional-probability}
	For events $D,H$ with $\Prob(H)>0$,
	\begin{equation}\label{eq:conditional-probability}
		\Prob(D\mid H)\eqdef\frac{\Prob(D\cap H)}{\Prob(H)} .
	\end{equation}
	Events $D$ and $H$ are \dfnas{independent events}{independent} if
	$\Prob(D\cap H)=\Prob(D)\Prob(H)$, equivalently
	$\Prob(D\mid H)=\Prob(D)$ when $\Prob(H)>0$.
\end{definition}

For fixed $H$ with $\Prob(H)>0$ the map $E\mapsto\Prob(E\mid H)$ is itself a
probability measure on $(\Omega,\mathcal{F})$: it assigns $1$ to $\Omega$ and is
$\sigma$-additive, both immediately from \eqref{eq:conditional-probability} and
$\sigma$-additivity of $\Prob$. Everything proved for probability measures
therefore applies to conditional probabilities without further argument.

\begin{theorem}[Bayes' rule]\label{thm:bayes}
	Let $D,H$ be events with $\Prob(D)>0$ and $\Prob(H)>0$. Then
	\begin{equation}\label{eq:bayes-1}
		\Prob(H\mid D)=\frac{\Prob(H)\,\Prob(D\mid H)}{\Prob(D)} .
	\end{equation}
	If moreover $0<\Prob(H)<1$,
	\begin{equation}\label{eq:bayes-2}
		\Prob(H\mid D)
		=\frac{\Prob(H)\,\Prob(D\mid H)}
		{\Prob(H)\,\Prob(D\mid H)+\Prob(H^{c})\,\Prob(D\mid H^{c})} .
	\end{equation}
	More generally, if $(H_j)_{j=1}^{n}$ partitions $\Omega$ with $\Prob(H_j)>0$ for
	every $j$, then for each $i$
	\begin{equation}\label{eq:bayes-partition}
		\Prob(H_i\mid D)
		=\frac{\Prob(D\mid H_i)\Prob(H_i)}
		{\sum_{j=1}^{n}\Prob(D\mid H_j)\Prob(H_j)} .
	\end{equation}
\end{theorem}

\begin{proof}
	By \eqref{eq:conditional-probability} applied twice,
	$\Prob(H\mid D)\Prob(D)=\Prob(D\cap H)=\Prob(D\mid H)\Prob(H)$; dividing by
	$\Prob(D)>0$ gives \eqref{eq:bayes-1}. For the denominators, the events
	$D\cap H_j$ are pairwise disjoint with union $D$ because $(H_j)$ partitions
	$\Omega$, so finite additivity (\autoref{prop:probability-basics}(iii)) gives
	the \dfnas{law of total probability}{law of total probability}
	\begin{equation}\label{eq:total-probability}
		\Prob(D)=\sum_{j=1}^{n}\Prob(D\cap H_j)
		=\sum_{j=1}^{n}\Prob(D\mid H_j)\Prob(H_j) .
	\end{equation}
	Substituting \eqref{eq:total-probability} into \eqref{eq:bayes-1} gives
	\eqref{eq:bayes-partition}, and \eqref{eq:bayes-2} is the case
	$n=2$ with $H_1=H$, $H_2=H^{c}$.
\end{proof}

The content of \eqref{eq:bayes-1} is a reversal of the direction of conditioning.
The likelihood $\Prob(D\mid H)$ runs from hypothesis to data, from cause to
effect; the posterior $\Prob(H\mid D)$ runs from data back to hypothesis, and
the theorem says the two are related by the ratio of the prior to the marginal
likelihood of the data \autocite[slide 6]{fukuda2026probability}.

\begin{eg}[Base rates dominate]\label{eg:medical-test}
	A disease affects $1\%$ of the population. The test is positive with probability
	$0.95$ on the diseased and negative with probability $0.9$ on the healthy. Let
	$H$ be ``diseased'' and $D$ ``tests positive''. Then
	\[
		\Prob(H\mid D)
		=\frac{0.01\times0.95}{0.01\times0.95+0.99\times0.10}
		=\frac{0.0095}{0.1085}=\frac{19}{217}\approx0.0876 .
	\]
	A test with $95\%$ sensitivity and $90\%$ specificity leaves the posterior
	probability of disease below $9\%$. The reason is arithmetic rather than
	statistical: the $99\%$ healthy generate $9.9$ false positives per hundred while
	the $1\%$ diseased generate $0.95$ true positives, and no reading of the
	likelihoods alone can reveal this. Neglect of the denominator in
	\eqref{eq:bayes-2} is the base-rate fallacy, and it is the mechanism behind
	overconfidence in models of belief updating.
\end{eg}

\begin{eg}[The airport carousel]\label{eg:carousel}
	Bags are correctly loaded with probability $0.9$. Conditional on the bag having
	been loaded, it fails to appear after nine minutes with probability $0.1$; if it
	was not loaded it certainly does not appear. With $H$ ``not loaded'' and $D$
	``bag missing'',
	\[
		\Prob(H\mid D)=\frac{0.1\times1}{0.1\times1+0.9\times0.1}
		=\frac{0.1}{0.19}=\frac{10}{19}\approx0.5263 .
	\]
	Waiting turns a prior of $0.1$ into a posterior above one half. The magnitude of
	the revision is governed by the likelihood ratio
	$\Prob(D\mid H)/\Prob(D\mid H^{c})=10$, which is the sufficient statistic
	for the evidence: \eqref{eq:bayes-2} in odds form reads
	\begin{equation}\label{eq:odds-form}
		\frac{\Prob(H\mid D)}{\Prob(H^{c}\mid D)}
		=\frac{\Prob(H)}{\Prob(H^{c})}\cdot
		\frac{\Prob(D\mid H)}{\Prob(D\mid H^{c})},
	\end{equation}
	posterior odds equal prior odds times the likelihood ratio.
\end{eg}

\begin{eg}[A poll]\label{eg:poll}
	Party $A$ has $30\%$ support, party $B$ has $40\%$, and $30\%$ support neither.
	The cabinet is supported by $80\%$ of $A$'s supporters, $20\%$ of $B$'s and
	$50\%$ of the rest. By \eqref{eq:bayes-partition} with the three-element
	partition, a randomly chosen cabinet supporter belongs to party $A$ with
	probability
	\[
		\frac{0.3\times0.8}{0.3\times0.8+0.4\times0.2+0.3\times0.5}
		=\frac{0.24}{0.47}=\frac{24}{47}\approx0.5106 .
	\]
\end{eg}

\subsection{Economic application: order flow and the market maker's beliefs}
\label{subsec:prob-microstructure}

\paragraph{Environment.} An asset has value $V\in\{H,L\}$ with $H>L$. A market
maker holds the prior $p\eqdef\Prob(V=H)\in(0,1)$ and quotes prices at which
they will trade one unit. A trader arrives; with probability $q\in(0,1)$ the
trader is \emph{informed} and knows $V$, and with probability $1-q$ the trader is
\emph{uninformed}. The informed trader buys when $V=H$ and sells when $V=L$; the
uninformed trader buys or sells with probability $\tfrac12$ each, independently of
$V$.

\begin{proposition}[Posterior after a sale]\label{prop:microstructure}
	Let $S$ denote the event that the first trade is a sale. Then
	\begin{equation}\label{eq:microstructure-posterior}
		\Prob(V=H\mid S)=\frac{p(1-q)}{p(1-q)+(1-p)(1+q)},
	\end{equation}
	which is strictly less than $p$ for every $q\in(0,1)$, is strictly decreasing in
	$q$, equals $p$ at $q=0$ and tends to $0$ as $q\to1$.
\end{proposition}

\begin{proof}
	Conditioning on the trader's type and using
	\eqref{eq:total-probability} within each state,
	\[
		\Prob(S\mid V=H)=q\cdot0+(1-q)\cdot\tfrac12=\tfrac{1-q}{2},
		\qquad
		\Prob(S\mid V=L)=q\cdot1+(1-q)\cdot\tfrac12=\tfrac{1+q}{2},
	\]
	because an informed trader never sells when the value is high and always sells
	when it is low. Substituting into \eqref{eq:bayes-2} and cancelling the factor
	$\tfrac12$ gives \eqref{eq:microstructure-posterior}. Writing the posterior as
	$\bigl(1+\frac{1-p}{p}\cdot\frac{1+q}{1-q}\bigr)^{-1}$ exhibits it as strictly
	decreasing in $q$, since $q\mapsto(1+q)/(1-q)$ is strictly increasing on
	$(0,1)$; the boundary values follow by inspection.
\end{proof}

The comparative static in $q$ is the informational content of order flow. When
$q=0$ every trade is noise and prices do not move; as $q\to1$ a single sale
reveals the state. The market maker who sets the bid at
$\E[V\mid S]=L+(H-L)\Prob(V=H\mid S)$ and the ask at $\E[V\mid B]$ makes
zero expected profit on each trade and covers the losses to informed traders with
the gains from uninformed ones; the difference between the two quotes is the
bid--ask spread, and \eqref{eq:microstructure-posterior} shows it widens in $q$.
This is the adverse-selection account of the spread, and the calculation above is
its entire mechanism \autocite[slide 12]{fukuda2026probability}; see
\textcite[Sec.~3.A.1]{ohara1998market} for the sequential-trade model built on
it and \textcite{kyle1985continuous} for the continuous version in which the
informed trader chooses order size strategically.

\section{Conditional Expectation}\label{sec:prob-conditional}

\begin{definition}[Conditional expectation given an event]\label{def:conditional-expectation-event}
	For an event $H$ with $\Prob(H)>0$ and a random variable $X$,
	\begin{equation}\label{eq:conditional-expectation}
		\E[X\mid H]\eqdef\sum_{x\in X(H)}x\,\Prob\bigl([X=x]\mid H\bigr)
		=\sum_{\omega\in H}X(\omega)\,\Prob\bigl(\{\omega\}\mid H\bigr)
		=\frac{\E[X\one_H]}{\Prob(H)},
	\end{equation}
	where $X(H)\eqdef\{X(\omega) : \omega\in H\}$ and $\one_H$ is the indicator of
	$H$.
\end{definition}

\begin{proof}[Proof of the equalities in \eqref{eq:conditional-expectation}]
	The first equals the second by grouping $\omega\in H$ according to the value of
	$X(\omega)$, exactly as in \eqref{eq:expectation}. For the third, expand
	$\Prob(\{\omega\}\mid H)=\Prob(\{\omega\})/\Prob(H)$ for $\omega\in H$ and
	note that $X(\omega)\one_H(\omega)=0$ for $\omega\notin H$, so
	$\sum_{\omega\in H}X(\omega)\Prob(\{\omega\})
		=\sum_{\omega\in\Omega}X(\omega)\one_H(\omega)\Prob(\{\omega\})=\E[X\one_H]$.
\end{proof}

\begin{definition}[Conditional expectation given a random variable]
	\label{def:conditional-expectation-rv}
	Let $X,Y$ be random variables. The \dfnas{conditional
		expectation}{conditional expectation} $\E[Y\mid X]$ is the random variable
	$\omega\mapsto\E\bigl[Y\mid[X=X(\omega)]\bigr]$, that is, the function of $X$
	whose value at $x\in X(\Omega)$ is
	\begin{equation}\label{eq:conditional-expectation-rv}
		\E[Y\mid X=x]\eqdef\sum_{y\in Y(\Omega)}y\,f_{Y\mid X}(y\mid x),
		\qquad
		f_{Y\mid X}(y\mid x)\eqdef\Prob\bigl([Y=y]\mid[X=x]\bigr).
	\end{equation}
	\nidx{condexpectation}{$\E[Y\mid X]$, conditional expectation}
\end{definition}

The distinction between \autoref{def:conditional-expectation-event} and
\autoref{def:conditional-expectation-rv} is that the first is a number and the
second a random variable --- a function of $X$, hence itself measurable with
respect to the information generated by $X$. Everything that follows turns on
this: $\E[Y\mid X]$ is the best summary of $Y$ available to someone who
observes $X$.

\begin{proposition}[Taking out what is known]\label{prop:taking-out}
	Let $X,Y$ be random variables with $Y$ integrable, and let
	$g\colon\R\to\R$ be such that $g(X)Y$ is integrable. Then
	\begin{equation}\label{eq:taking-out}
		\E\bigl[g(X)Y\mid X\bigr]=g(X)\,\E[Y\mid X] .
	\end{equation}
	In particular $\E[g(X)\mid X]=g(X)$ and $\E[a\mid X]=a$ for constant $a$.
\end{proposition}

\begin{proof}
	Fix $x_0\in X(\Omega)$ and put $W\eqdef g(X)Y$. If $g(x_0)=0$ then
	$W$ is identically zero on $[X=x_0]$ and both sides of \eqref{eq:taking-out}
	vanish there. If $g(x_0)\neq0$ then on the event $[X=x_0]$ we have
	$W=w$ if and only if $Y=w/g(x_0)$, so
	\[
		\Prob\bigl([W=w]\mid[X=x_0]\bigr)
		=\Prob\Bigl(\bigl[Y=\tfrac{w}{g(x_0)}\bigr]\mid[X=x_0]\Bigr),
	\]
	and substituting $y=w/g(x_0)$ in \eqref{eq:conditional-expectation-rv},
	\[
		\E[W\mid X=x_0]=\sum_{w}w\,\Prob\bigl([W=w]\mid[X=x_0]\bigr)
		=g(x_0)\sum_{y}y\,\Prob\bigl([Y=y]\mid[X=x_0]\bigr)
		=g(x_0)\E[Y\mid X=x_0].
	\]
	Since $x_0$ was arbitrary, \eqref{eq:taking-out} holds as an identity between
	random variables. Taking $Y\equiv1$ gives $\E[g(X)\mid X]=g(X)$, and $g$
	constant gives the last claim.
\end{proof}

Both integrability hypotheses are needed, and neither implies the other. With
$g\equiv0$ the product $g(X)Y$ is integrable whatever $Y$ is, so the left side
of \eqref{eq:taking-out} is $0$, while the right side contains the factor
$\E[Y\mid X]$, which is undefined for non-integrable $Y$. Conversely, an
integrable $Y$ and an unbounded $g$ can make $g(X)Y$ non-integrable.

\begin{theorem}[Law of iterated expectations]\label{thm:lie}
	\begin{enumerate}[(i)]
		\item Let $(H_j)_{j=1}^{m}$ partition $\Omega$ with $\Prob(H_j)>0$ for every
		      $j$, and let $X$ be integrable. Then
		      \begin{equation}\label{eq:lie-partition}
			      \E[X]=\sum_{j=1}^{m}\E[X\mid H_j]\,\Prob(H_j).
		      \end{equation}
		\item For integrable $X$ and any random variable $Y$,
		      \begin{equation}\label{eq:lie}
			      \E[X]=\E\bigl[\E[X\mid Y]\bigr].
		      \end{equation}
	\end{enumerate}
\end{theorem}

\begin{proof}
	(i) By the third form of \eqref{eq:conditional-expectation},
	$\E[X\mid H_j]\Prob(H_j)=\E[X\one_{H_j}]$. Since $(H_j)$ partitions $\Omega$,
	$\sum_j\one_{H_j}=1$ pointwise, so linearity
	(\autoref{prop:expectation-properties}(i)) gives
	$\sum_j\E[X\one_{H_j}]=\E\bigl[X\sum_j\one_{H_j}\bigr]=\E[X]$.

	(ii) Apply (i) to the partition $H_j\eqdef[Y=y_j]$, $y_j\in Y(\Omega)$, and
	observe that $\sum_j\E[X\mid Y=y_j]\Prob(Y=y_j)$ is by
	\eqref{eq:expectation} and \eqref{eq:conditional-expectation-rv} precisely the
	expectation of the random variable $\E[X\mid Y]$, whose value on $[Y=y_j]$ is
	$\E[X\mid Y=y_j]$ and which takes that value with probability
	$\Prob(Y=y_j)$.
\end{proof}

These are Theorems 2.4 and 2.5 of \textcite[pp.~74--75]{fukuda2026notes}
\autocite[slides 18--19]{fukuda2026probability}. The statement to hold on to is
that conditioning on less information cannot change the average: an expectation
computed in two stages, first within each cell of a partition and then across
cells, equals the expectation computed at once. Almost every argument in dynamic
economics that replaces a many-period expectation by a one-period expectation of
a continuation value is an application of \eqref{eq:lie}.

\begin{proposition}[Conditional Jensen]\label{prop:conditional-jensen}
	Let $\varphi\colon I\to\R$ be convex on an interval $I$, and let $Y$ be
	integrable with $\Prob(Y\in I)=1$ and $\E[Y\mid X]$ almost surely interior to
	$I$. Then
	\begin{equation}\label{eq:conditional-jensen}
		\E\bigl[\varphi(Y)\mid X\bigr]\geq\varphi\bigl(\E[Y\mid X]\bigr)
		\qquad\text{almost surely,}
	\end{equation}
	and the inequality reverses for concave $\varphi$.
\end{proposition}

\begin{proof}
	Fix $x\in X(\Omega)$. Conditional on $[X=x]$ the map
	$E\mapsto\Prob(E\mid[X=x])$ is a probability measure, as noted after
	\autoref{def:conditional-probability}, and $\E[\cdot\mid X=x]$ is the
	expectation under it. \autoref{thm:jensen} applied under that measure, with
	$\mu\eqdef\E[Y\mid X=x]$ interior to $I$, gives
	$\E[\varphi(Y)\mid X=x]\geq\varphi(\mu)$. As $x$ was arbitrary, the inequality
	holds pointwise as an inequality between random variables.
\end{proof}

Taking expectations in \eqref{eq:conditional-jensen} and applying
\eqref{eq:lie} recovers \autoref{thm:jensen}, so the conditional version is
strictly stronger. It is what makes the value function of a dynamic programme
concave in the state when the return function is concave, and it is what forces
the martingale $\bigl(\E[Y\mid X_t]\bigr)_t$ to have a convex-increasing
transform.

\subsection{Conditioning on a \texorpdfstring{$\sigma$}{sigma}-algebra}
\label{subsec:prob-general-conditioning}

Everything above conditions on an event of positive probability. In the models of
\autoref{sec:prob-normal} and \autoref{ch:dynamic} the conditioning variable is
continuously distributed, so $\Prob([X=x])=0$ for every $x$ and
\autoref{def:conditional-probability} does not apply. The construction that
replaces it takes the conditioning \emph{information} rather than a conditioning
event as primitive.

\begin{definition}[Conditional expectation given a $\sigma$-algebra]
	\label{def:conditional-expectation-sigma}
	Let $(\Omega,\mathcal{F},\Prob)$ be a probability space, $\mathcal{G}\subseteq\mathcal{F}$ a
	$\sigma$-algebra, and $Y$ an integrable random variable. A random variable $Z$
	is a \dfnas{conditional expectation!given a
		$\sigma$-algebra}{conditional expectation} of $Y$ given $\mathcal{G}$, written $Z=\E[Y\mid\mathcal{G}]$, if
	\begin{enumerate}[(i)]
		\item $Z$ is $\mathcal{G}$-measurable and integrable; and
		\item $\E[Z\one_G]=\E[Y\one_G]$ for every $G\in\mathcal{G}$.
	\end{enumerate}
	\nidx{condexpsigma}{$\E[Y\mid\mathcal{G}]$, conditional expectation given $\mathcal{G}$}
\end{definition}

Condition (i) is the requirement that the forecast use only the information in
$\mathcal{G}$; condition (ii) is the requirement that it be unbiased on every event that
information can distinguish. Nothing here refers to the probability of a single
value of a conditioning variable.

\begin{theorem}[Existence and almost sure uniqueness]\label{thm:cond-exp-exists}
	Let $\mathcal{G}\subseteq\mathcal{F}$ be a $\sigma$-algebra and let $Y$ be integrable. Then
	$\E[Y\mid\mathcal{G}]$ exists, and any two versions agree almost surely. If in addition
	$Y\in L^{2}$, then $\E[Y\mid\mathcal{G}]$ is the orthogonal projection of $Y$ onto the
	closed subspace $L^{2}(\Omega,\mathcal{G},\Prob)\subseteq L^{2}(\Omega,\mathcal{F},\Prob)$.
\end{theorem}

\begin{proof}
	\emph{Uniqueness.} Let $Z,Z'$ both satisfy (i) and (ii). The event
	$G\eqdef[Z>Z']$ lies in $\mathcal{G}$ because both variables are $\mathcal{G}$-measurable, and
	(ii) applied to each gives $\E[(Z-Z')\one_G]=\E[Y\one_G]-\E[Y\one_G]=0$. The
	integrand is strictly positive on $G$ and zero off it, so
	\autoref{prop:expectation-properties}(ii) gives $\Prob(G)=0$. Exchanging the
	roles of $Z$ and $Z'$ gives $\Prob([Z'>Z])=0$, hence $Z=Z'$ almost surely.

	\emph{Existence for $Y\in L^{2}$.} The set
	$L^{2}(\Omega,\mathcal{G},\Prob)$ of square-integrable $\mathcal{G}$-measurable random variables
	is a subspace of $L^{2}(\Omega,\mathcal{F},\Prob)$, and it is closed: an $L^{2}$-convergent
	sequence has an almost surely convergent subsequence, whose limit is again
	$\mathcal{G}$-measurable as a pointwise limit of $\mathcal{G}$-measurable functions. Let
	$Z\eqdef P_{L^{2}(\mathcal{G})}Y$, which exists and is unique by
	\autoref{thm:orthogonal-decomposition}. For $G\in\mathcal{G}$ the variable $\one_G$ lies
	in $L^{2}(\Omega,\mathcal{G},\Prob)$, so \autoref{prop:proj-subspace-char} gives
	$\E[(Y-Z)\one_G]=\ip{Y-Z}{\one_G}=0$, which is (ii); and $Z$ is
	$\mathcal{G}$-measurable by construction, which is (i).

	\emph{Existence for integrable $Y$.} The extension from $L^{2}$ to $L^{1}$ is a
	monotone-convergence argument that this volume does not develop; the direct
	construction applies the Radon--Nikodym theorem to the finite measure
	$G\mapsto\E[Y^{+}\one_G]$ on $(\Omega,\mathcal{G})$, which is absolutely continuous with
	respect to the restriction of $\Prob$, and subtracts the corresponding object
	for $Y^{-}$. The statement and this proof are Theorem~11.3 of
	\textcite[p.~230]{legall2022measure}, resting on the Radon--Nikodym theorem
	\autocite[Thm.~4.11, p.~75]{legall2022measure}.
\end{proof}

Two further facts complete the framework, and both are used without comment in
applied work.

\begin{theorem}[Doob--Dynkin]\label{thm:doob-dynkin}
	Let $X$ be a random variable and $\sigma(X)\eqdef\{X^{-1}(B) : B\
		\text{Borel}\}$ the $\sigma$-algebra it generates. A random variable $Z$ is
	$\sigma(X)$-measurable if and only if $Z=\varphi(X)$ for some Borel
	$\varphi\colon\R\to\R$.
\end{theorem}

Consequently $\E[Y\mid X]\eqdef\E[Y\mid\sigma(X)]$ can be written as $\varphi(X)$,
and the symbol $\E[Y\mid X=x]$ is \emph{defined} to mean $\varphi(x)$. The
function $\varphi$ is determined only up to sets of $\Prob_X$-measure zero, so
$\E[Y\mid X=x]$ has no meaning at an individual $x$ in isolation; statements about
it are statements holding for $\Prob_X$-almost every $x$. This is the precise
sense in which conditioning on a null event is legitimate.

\begin{proposition}[Evaluation with densities]\label{prop:cond-exp-density}
	Let $(X,Y)$ have joint density $p$ on $\R^2$ and let $q(x)=\int_\R p(x,y)\diff y$
	be the marginal density of $X$. If $Y$ is integrable, then a version of
	$\E[Y\mid X]$ is $\varphi(X)$ with
	\begin{equation}\label{eq:cond-exp-density}
		\varphi(x)=\frac{1}{q(x)}\int_{\R}y\,p(x,y)\diff y
		\qquad\text{for every }x\text{ with }q(x)>0,
	\end{equation}
	and $\varphi(x)$ arbitrary where $q(x)=0$; the exceptional set has
	$\Prob_X$-measure zero.
\end{proposition}

This is Proposition~11.11 of \textcite[p.~244]{legall2022measure}, stated there
for non-negative $h(X)$ and general dimensions. Formula
\eqref{eq:cond-exp-density} is the one used in \autoref{sec:prob-normal}, and it
is the exact analogue of \eqref{eq:conditional-expectation-rv} with the
conditional mass function $f_{Y\mid X}$ replaced by the conditional density
$p(x,\cdot)/q(x)$.

\begin{proposition}[The discrete case is a special case]\label{prop:cond-exp-discrete}
	Let $X$ take countably many values and let $Y$ be integrable. Then
	$\E[Y\mid\sigma(X)]=\varphi(X)$ almost surely, where
	$\varphi(x)=\E[Y\mid[X=x]]$ in the sense of
	\autoref{def:conditional-expectation-event} for every $x$ with $\Prob([X=x])>0$.
\end{proposition}

\begin{proof}
	The atoms of $\sigma(X)$ are the events $[X=x]$ with $\Prob([X=x])>0$, together
	with a null set, so a $\sigma(X)$-measurable variable is constant on each atom
	and $Z\eqdef\varphi(X)$ satisfies (i). For (ii) it suffices, by countable
	additivity and $\sigma$-additivity of the expectation, to check
	$\E[Z\one_G]=\E[Y\one_G]$ on the atoms $G=[X=x]$, where both sides equal
	$\E[Y\one_{[X=x]}]$ by the third form of \eqref{eq:conditional-expectation}.
\end{proof}

Every property established above in the countable setting ---
\autoref{prop:taking-out}, \autoref{thm:lie} and
\autoref{prop:conditional-jensen} --- holds for $\E[\cdot\mid\mathcal{G}]$ with ``almost
surely'' appended, and the tower property takes the sharper form
$\E\bigl[\E[Y\mid\mathcal{F}_2]\mid\mathcal{F}_1\bigr]=\E[Y\mid\mathcal{F}_1]$ for
$\mathcal{F}_1\subseteq\mathcal{F}_2$: coarser information wins. The proofs are those of
\textcite[pp.~234--235, 238--241]{legall2022measure}. In the $L^{2}$ case they are
also immediate from \autoref{thm:cond-exp-exists}, since composing orthogonal
projections onto nested closed subspaces gives the projection onto the smaller
one --- which is the content of \autoref{ex:proj-1} and the reason
\autoref{sec:prob-projection} can treat conditioning as geometry.

\section{Conditional Expectation as a Projection}\label{sec:prob-projection}

Let $L^{2}\eqdef L^{2}(\Omega,\mathcal{F},\Prob)$ be the space of square-integrable random
variables with the inner product $\ip{U}{V}\eqdef\E[UV]$; it is a Hilbert space,
and the geometry of \autoref{ch:projection} applies verbatim with
$\norm{U}^{2}=\E[U^{2}]$.

\begin{theorem}[Conditional expectation is the best predictor]\label{thm:best-predictor}
	Let $X,Y\in L^{2}$ and let
	$M(X)\eqdef\{\phi(X) : \phi\ \text{Borel},\ \phi(X)\in L^{2}\}$, a closed
	subspace of $L^{2}$. Then for every $\phi$ with $\phi(X)\in L^{2}$,
	\begin{equation}\label{eq:best-predictor}
		\E\bigl[(Y-\phi(X))^{2}\bigr]
		=\E\bigl[(Y-\E[Y\mid X])^{2}\bigr]
		+\E\bigl[(\E[Y\mid X]-\phi(X))^{2}\bigr]
		\geq\E\bigl[(Y-\E[Y\mid X])^{2}\bigr],
	\end{equation}
	with equality if and only if $\phi(X)=\E[Y\mid X]$ almost surely. Equivalently,
	$\E[Y\mid X]$ is the orthogonal projection of $Y$ onto $M(X)$:
	$P_{M(X)}Y=\E[Y\mid X]$.
\end{theorem}

\begin{proof}
	Write $\hat Y\eqdef\E[Y\mid X]$ and decompose
	$Y-\phi(X)=(Y-\hat Y)+(\hat Y-\phi(X))$. Expanding the square,
	\[
		\E\bigl[(Y-\phi(X))^{2}\bigr]
		=\E\bigl[(Y-\hat Y)^{2}\bigr]
		+2\,\E\bigl[(Y-\hat Y)(\hat Y-\phi(X))\bigr]
		+\E\bigl[(\hat Y-\phi(X))^{2}\bigr].
	\]
	The cross term vanishes. Indeed $\hat Y-\phi(X)$ is a function of $X$, so
	\autoref{prop:taking-out} and then \eqref{eq:lie} give
	\[
		\E\bigl[(Y-\hat Y)(\hat Y-\phi(X))\bigr]
		=\E\Bigl[\E\bigl[(Y-\hat Y)(\hat Y-\phi(X))\mid X\bigr]\Bigr]
		=\E\Bigl[(\hat Y-\phi(X))\,\underbrace{\bigl(\E[Y\mid X]-\hat Y\bigr)}_{=0}\Bigr]
		=0,
	\]
	using $\E[\hat Y\mid X]=\hat Y$ from \autoref{prop:taking-out}. This proves the
	identity in \eqref{eq:best-predictor}; the inequality and its equality case
	follow because the second term is non-negative and vanishes only when
	$\hat Y=\phi(X)$ almost surely, by
	\autoref{prop:expectation-properties}(ii).

	The vanishing of the cross term for \emph{every} $\phi$ says
	$Y-\hat Y\perp M(X)$, and $\hat Y\in M(X)$; by
	\autoref{prop:proj-subspace-char} this characterises $\hat Y$ as the orthogonal
	projection of $Y$ onto $M(X)$, whose existence and uniqueness are
	\autoref{thm:orthogonal-decomposition}.
\end{proof}

Equation \eqref{eq:best-predictor} carries three readings. As a minimisation, it says the conditional expectation is the mean-square-optimal
forecast among all functions of the conditioning variable, extending the
constant-forecast statement of \autoref{prop:variance-identity}. As a geometric
statement, it identifies conditioning with orthogonal projection, so that all the
structure of \autoref{ch:projection} --- idempotence, the Pythagorean
decomposition, contraction in norm --- transfers to conditioning; the tower
property \eqref{eq:lie} becomes the composition of projections onto nested
subspaces. As a variance decomposition, taking $\phi\equiv\E[Y]$ in
\eqref{eq:best-predictor} and using \eqref{eq:lie} yields
\begin{equation}\label{eq:variance-decomposition}
	\Var(Y)=\E\bigl[\Var(Y\mid X)\bigr]+\Var\bigl(\E[Y\mid X]\bigr),
\end{equation}
the law of total variance: unexplained variation plus explained variation. In a
regression context the second term over the first is the population $R^{2}$.

\begin{remark}[Linear versus general prediction]\label{rem:linear-vs-general}
	\autoref{prop:best-linear-predictor} solved the same problem over the strictly
	smaller subspace $\Span\{1,X\}\subseteq M(X)$ and produced
	\begin{equation}\label{eq:blp-again}
		\E^{\ast}[Y\mid X]=\E[Y]+\frac{\Cov(X,Y)}{\Var(X)}\bigl(X-\E[X]\bigr).
	\end{equation}
	Since the subspace is smaller, the linear predictor is weakly worse, and
	strictly worse whenever $\E[Y\mid X]$ is not affine in $X$. The two coincide
	exactly when $\E[Y\mid X]$ happens to be affine, which is the content of the
	next proposition and the reason the bivariate normal is the workhorse of applied
	work: it is the case in which linear regression estimates the conditional
	expectation rather than merely its best affine approximation.
\end{remark}

\begin{proposition}[The Gaussian case]\label{prop:bivariate-normal}
	Let $(X,Y)$ be jointly normal with $\Var(X)>0$. Then
	\begin{equation}\label{eq:gaussian-conditional}
		\E[Y\mid X]=\E[Y]+\frac{\Cov(X,Y)}{\Var(X)}\bigl(X-\E[X]\bigr),
		\qquad
		\Var(Y\mid X)=\bigl(1-\Corr(X,Y)^{2}\bigr)\Var(Y),
	\end{equation}
	the conditional variance being a constant, not a function of $X$.
\end{proposition}

The multivariate statement, in which $X$ and $Y$ are vectors and the ratio
$\Cov(X,Y)/\Var(X)$ becomes $\Sigma_{12}\Sigma_{22}^{-1}$, is
\autoref{thm:mvn-conditional} of \autoref{app:normal}, together with the affine
stability of the family that makes the argument work.

\begin{proof}
	Let $g(X)$ denote the right-hand side of \eqref{eq:blp-again}. By construction
	$Y-g(X)$ is orthogonal to $1$ and to $X$, that is $\E[Y-g(X)]=0$ and
	$\Cov\bigl(X,Y-g(X)\bigr)=0$. Now $Y-g(X)$ and $X$ are jointly normal, being
	affine images of the jointly normal pair $(X,Y)$, and for jointly normal
	variables zero correlation implies independence, since the joint density then
	factors. Hence
	$\E[Y-g(X)\mid X]=\E[Y-g(X)]=0$, which is the first identity in
	\eqref{eq:gaussian-conditional}.

	For the variance, $Y-\E[Y]=\frac{\Cov(X,Y)}{\Var(X)}(X-\E[X])+(Y-g(X))$ is an
	orthogonal decomposition, so \autoref{prop:pythagoras} gives
	\[
		\Var(Y)=\frac{\Cov(X,Y)^{2}}{\Var(X)^{2}}\Var(X)+\Var\bigl(Y-g(X)\bigr)
		=\frac{\Cov(X,Y)^{2}}{\Var(X)}+\Var(Y\mid X),
	\]
	and rearranging gives the stated expression with
	$\Corr(X,Y)^{2}=\Cov(X,Y)^{2}/(\Var(X)\Var(Y))$.
\end{proof}

\section{Normal--Normal Updating}\label{sec:prob-normal}

\begin{theorem}[Gaussian updating, one observation]\label{thm:normal-updating}
	Let $X\sim\mathcal{N}(\mu_x,\sigma_x^{2})$ with $\sigma_x^2>0$ and let
	$Y=X+\varepsilon$ with $\varepsilon\sim\mathcal{N}(0,\sigma_\varepsilon^{2})$,
	$\sigma_\varepsilon^2>0$, independent of $X$. Then
	\begin{equation}\label{eq:normal-updating}
		X\mid Y=y\ \sim\
		\mathcal{N}\!\left(
		\frac{\sigma_\varepsilon^{2}\mu_x+\sigma_x^{2}y}
		{\sigma_x^{2}+\sigma_\varepsilon^{2}},\
		\frac{\sigma_x^{2}\sigma_\varepsilon^{2}}
		{\sigma_x^{2}+\sigma_\varepsilon^{2}}
		\right).
	\end{equation}
	Writing $\tau_x\eqdef1/\sigma_x^{2}$ and
	$\tau_\varepsilon\eqdef1/\sigma_\varepsilon^{2}$ for the \dfnas{precision}{precisions}
	and $\tau\eqdef\tau_x+\tau_\varepsilon$, this reads
	\begin{equation}\label{eq:normal-updating-precision}
		X\mid Y=y\ \sim\
		\mathcal{N}\!\left(\frac{\tau_x\mu_x+\tau_\varepsilon y}{\tau},\ \tau^{-1}\right).
	\end{equation}
\end{theorem}

\begin{proof}
	The prior density and the likelihood are
	\[
		\pi(x)=\frac{1}{\sqrt{2\pi}\sigma_x}
		\exp\!\left(-\frac{(x-\mu_x)^{2}}{2\sigma_x^{2}}\right),
		\qquad
		L(y\mid x)=\frac{1}{\sqrt{2\pi}\sigma_\varepsilon}
		\exp\!\left(-\frac{(y-x)^{2}}{2\sigma_\varepsilon^{2}}\right),
	\]
	the second because $Y-X=\varepsilon$ is independent of $X$. By Bayes' rule in
	density form,
	$f(x\mid Y=y)=\pi(x)L(y\mid x)\big/\int_{-\infty}^{\infty}
		\pi(\tilde x)L(y\mid\tilde x)\,\diff\tilde x$, so as a function of $x$,
	\[
		f(x\mid Y=y)\propto
		\exp\!\left(-\frac{(x-\mu_x)^{2}}{2\sigma_x^{2}}
		-\frac{(y-x)^{2}}{2\sigma_\varepsilon^{2}}\right).
	\]
	Complete the square in the exponent. Multiplying out,
	\begin{align*}
		-\frac{(x-\mu_x)^{2}}{2\sigma_x^{2}}-\frac{(y-x)^{2}}{2\sigma_\varepsilon^{2}}
		 & =-\frac{1}{2\sigma_x^{2}\sigma_\varepsilon^{2}}
		\Bigl[\sigma_\varepsilon^{2}(x-\mu_x)^{2}+\sigma_x^{2}(y-x)^{2}\Bigr]         \\
		 & =-\frac{1}{2\sigma_x^{2}\sigma_\varepsilon^{2}}
		\Bigl[(\sigma_x^{2}+\sigma_\varepsilon^{2})x^{2}
		-2(\sigma_\varepsilon^{2}\mu_x+\sigma_x^{2}y)x
		+(\sigma_\varepsilon^{2}\mu_x^{2}+\sigma_x^{2}y^{2})\Bigr]                    \\
		 & =-\frac{\sigma_x^{2}+\sigma_\varepsilon^{2}}
		{2\sigma_x^{2}\sigma_\varepsilon^{2}}
		\left(x-\frac{\sigma_\varepsilon^{2}\mu_x+\sigma_x^{2}y}
		{\sigma_x^{2}+\sigma_\varepsilon^{2}}\right)^{2}+c(y),
	\end{align*}
	where $c(y)$ collects all terms free of $x$. Exponentiating, the posterior
	density is proportional to a Gaussian kernel in $x$ with mean
	$(\sigma_\varepsilon^{2}\mu_x+\sigma_x^{2}y)/(\sigma_x^{2}+\sigma_\varepsilon^{2})$
	and variance
	$\sigma_x^{2}\sigma_\varepsilon^{2}/(\sigma_x^{2}+\sigma_\varepsilon^{2})$; since
	a density is determined by its kernel up to the normalising constant, and by
	\autoref{prop:gaussian-integral} that constant is
	$\bigl(2\pi\sigma_x^{2}\sigma_\varepsilon^{2}/
		(\sigma_x^{2}+\sigma_\varepsilon^{2})\bigr)^{-1/2}$,
	\eqref{eq:normal-updating} follows. The precision form is the substitution
	$\sigma_x^{2}=1/\tau_x$, $\sigma_\varepsilon^{2}=1/\tau_\varepsilon$, under which
	$\sigma_x^2+\sigma_\varepsilon^2=\tau/(\tau_x\tau_\varepsilon)$.
\end{proof}

This is the conjugate normal model of \textcite[Sec.~2.5, pp.~39--42]{gelman2013bayesian},
whose multivariate counterpart with known variance is Section~3.5 of the same
work; the precision form \eqref{eq:normal-updating-precision} is the statement
there that the posterior precision is the prior precision plus the data
precision. Three of its properties are the reason the Gaussian model is used
wherever information is aggregated.

\emph{Precisions add.} The posterior precision $\tau=\tau_x+\tau_\varepsilon$
does not depend on the realisation $y$: the informativeness of an observation is
known before it is made. Consequently the posterior variance
$\tau^{-1}<\sigma_x^{2}$ falls with every observation, deterministically.

\emph{The posterior mean is a precision-weighted average.} It is linear in $y$,
\[
	\E[X\mid Y=y]
	=\mu_x+\frac{\sigma_x^{2}}{\sigma_x^{2}+\sigma_\varepsilon^{2}}(y-\mu_x),
\]
so the response to the surprise $y-\mu_x$ is the ratio of prior variance to total
variance, between $0$ and $1$. A precise prior or a noisy signal produces
sluggish updating; the converse produces aggressive updating. Linearity in $y$ is
also the statement that the best predictor is the best \emph{linear} predictor,
which is \autoref{prop:bivariate-normal} applied to the jointly normal pair
$(X,Y)$.

\emph{Conjugacy.} The posterior is again normal, so the operation can be iterated
without leaving the family --- and this is what makes recursive filtering
tractable.

\begin{theorem}[Gaussian updating, many observations]\label{thm:normal-updating-many}
	Let $X\sim\mathcal{N}(\mu_x,\tau_x^{-1})$ and let
	$Y_j=X+\varepsilon_j$ for $j=1,\dots,n$, where the
	$\varepsilon_j\sim\mathcal{N}(0,\tau_j^{-1})$ are mutually independent and
	independent of $X$. Then with $\tau\eqdef\tau_x+\sum_{j=1}^{n}\tau_j$,
	\begin{equation}\label{eq:normal-updating-many}
		X\mid(Y_j=y_j)_{j=1}^{n}\ \sim\
		\mathcal{N}\!\left(\frac{\tau_x}{\tau}\mu_x
		+\sum_{j=1}^{n}\frac{\tau_j}{\tau}y_j,\ \tau^{-1}\right).
	\end{equation}
\end{theorem}

\begin{proof}
	Independence of the errors makes the joint likelihood the product
	$\prod_{j}L(y_j\mid x)$, so
	\[
		f\bigl(x\mid(Y_j=y_j)_j\bigr)\propto
		\exp\!\left(-\tfrac12\Bigl[\tau_x(x-\mu_x)^{2}
			+\sum_{j=1}^{n}\tau_j(x-y_j)^{2}\Bigr]\right).
	\]
	The bracket is a quadratic in $x$ with leading coefficient
	$\tau_x+\sum_j\tau_j=\tau$ and linear coefficient
	$-2\bigl(\tau_x\mu_x+\sum_j\tau_jy_j\bigr)$, so completing the square as in
	\autoref{thm:normal-updating} gives
	$\tau\bigl(x-(\tau_x\mu_x+\sum_j\tau_jy_j)/\tau\bigr)^{2}$ plus terms free of
	$x$. Identification of the Gaussian kernel gives
	\eqref{eq:normal-updating-many}.
\end{proof}

The posterior mean is a weighted average of the prior mean and all the
observations, with weights proportional to precisions summing to one, and the
posterior precision is the sum of all precisions. Applying
\autoref{thm:normal-updating} sequentially, using the posterior after $j$
observations as the prior for observation $j+1$, gives the same answer --- the
order of arrival is irrelevant, which is a substantive restriction and not an
accounting identity.

\subsection{Economic application: CARA--normal portfolio choice}
\label{subsec:prob-cara}

\paragraph{Environment.} An investor with initial wealth $W_0$ and constant
absolute risk aversion $a>0$, $u(W)=-e^{-aW}$, chooses a quantity $\theta\in\R$
of a risky asset with payoff $V\sim\mathcal{N}(\mu,\sigma^{2})$ at price $p$; the
riskless rate is normalised to zero, so terminal wealth is
$W=W_0+\theta(V-p)$.

\begin{proposition}[Gaussian demand]\label{prop:cara-demand}
	The investor's optimal holding is
	\begin{equation}\label{eq:cara-demand}
		\theta^{\ast}=\frac{\mu-p}{a\sigma^{2}},
	\end{equation}
	and the certainty equivalent of the optimal portfolio is
	$W_0+(\mu-p)^{2}/(2a\sigma^{2})$. If instead the investor observes a signal
	$Y=V+\varepsilon$ with $\varepsilon\sim\mathcal{N}(0,\tau_\varepsilon^{-1})$
	independent of $V$, then $\theta^{\ast}$ is \eqref{eq:cara-demand} with $\mu$
	and $\sigma^{2}$ replaced by the posterior mean and variance of
	\autoref{thm:normal-updating}.
\end{proposition}

\begin{proof}
	Terminal wealth is normal with mean $W_0+\theta(\mu-p)$ and variance
	$\theta^{2}\sigma^{2}$, so by \autoref{prop:exponential-moments}
	\[
		\E[u(W)]=-\exp\!\left(-a\bigl(W_0+\theta(\mu-p)\bigr)
		+\tfrac12a^{2}\theta^{2}\sigma^{2}\right).
	\]
	Maximising this is minimising the exponent, that is maximising
	$\theta(\mu-p)-\tfrac12a\theta^{2}\sigma^{2}$, a strictly concave quadratic in
	$\theta$; \autoref{thm:foc-concave} gives the unique maximiser
	\eqref{eq:cara-demand}, and substituting back gives an exponent of
	$-a\bigl(W_0+(\mu-p)^{2}/(2a\sigma^{2})\bigr)$, whose certainty equivalent is
	the stated value. Conditioning on $Y=y$ replaces the distribution of $V$ by its
	posterior, which is again normal by \autoref{thm:normal-updating}, and the same
	computation applies.
\end{proof}

Demand is linear in the price with slope $-1/(a\sigma^{2})$, and initial wealth
does not appear --- the wealth effect is exactly zero under CARA, which is what
makes the model aggregate. Two consequences follow. Because demand is linear in
the posterior mean and the posterior mean is linear in the signal by
\autoref{thm:normal-updating}, individual demands are linear in private signals,
so aggregate demand is linear in the average signal and the market-clearing price
is a linear statistic of the population's information. This is the mechanism by
which prices aggregate dispersed information; adding noise traders to prevent the
price from being fully revealing gives the noisy rational-expectations
equilibrium of \textcite{grossman1980impossibility}, in which the price is
informative but not sufficient, so that costly information acquisition remains
worthwhile. Second, the precision arithmetic of \eqref{eq:normal-updating-precision}
means an investor with a more precise signal takes a larger position on the same
price, which is the source of the informed trader's profit in
\autoref{subsec:prob-microstructure}.

\subsection{Economic application: recursive filtering}\label{subsec:prob-kalman}

Conjugacy makes updating recursive. Let the state follow
$X_{t+1}=\rho X_t+\eta_{t+1}$ with $\eta_{t+1}\sim\mathcal{N}(0,\tau_2)$ and let the
observation be $Y_t=X_t+\varepsilon_t$ with
$\varepsilon_t\sim\mathcal{N}(0,\tau_1^{-1})$, all errors independent. If the prior
on $X_t$ given the history is $\mathcal{N}(\hat x_t,P_t)$, then
\autoref{thm:normal-updating} gives the posterior precision
$P_t^{-1}+\tau_1$ after observing $Y_t$, and the transition inflates the variance
by $\tau_2$ after scaling by $\rho^{2}$:
\begin{equation}\label{eq:riccati-recursion}
	P_{t+1}=\rho^{2}\bigl(P_t^{-1}+\tau_1\bigr)^{-1}+\tau_2 .
\end{equation}
The right-hand side is exactly the operator $T$ of \eqref{eq:riccati}, shown in
\autoref{prop:riccati-contraction} to be a contraction of modulus $\rho^{2}$ on
$\R_{++}$. \autoref{thm:banach-fixed-point} therefore gives a unique steady-state
variance $x^{\ast}$ satisfying \eqref{eq:riccati-fixed-point}, reached from every
initial uncertainty at the geometric rate $\rho^{2}$. The economics of the
recursion is that the filter never learns the state exactly, no matter how long
it observes: each period's transition injects fresh variance $\tau_2$, and the
steady state balances that injection against the precision $\tau_1$ acquired from
the new observation. Setting $\rho=0$ makes the state i.i.d.\ and gives
$x^{\ast}=\tau_2$, the variance of a state about which yesterday's data say
nothing.

Equation \eqref{eq:riccati-recursion} is the scalar case of the Kalman
recursions of \textcite[Sec.~9.4, p.~270]{brockwell2016introduction}, whose
general form replaces the variances by covariance matrices and the scalar $\rho$
by a transition matrix $A$:
\[
	P_{t+1}=A\bigl(P_t^{-1}+C^{\transp}R^{-1}C\bigr)^{-1}A^{\transp}+Q,
\]
with $C$ the observation matrix, $R$ the observation-noise covariance and $Q$ the
state-noise covariance. The contraction argument does \emph{not} carry over
verbatim. The scalar proof uses that $x\mapsto\rho^2(x^{-1}+\tau_1)^{-1}+\tau_2$
is a contraction for the Euclidean metric on $\R_{++}$, and the matrix analogue
is false for the operator norm on positive definite matrices. What is true is
that the discrete algebraic Riccati equation has a unique positive semidefinite
stabilising solution, to which $P_t$ converges from every positive semidefinite
initialisation, provided the pair $(A,Q^{1/2})$ is stabilisable and the pair
$(A,C)$ is detectable; the natural argument uses a Riemannian metric on the cone
of positive definite matrices rather than the Euclidean one. Stabilisability and
detectability are the conditions under which unobserved directions of the state
decay on their own and observed directions are pinned down by the data; without
them the filter can carry an undamped, unmeasured mode and the error covariance
need not settle. See \textcite[Sec.~9.4, p.~270]{brockwell2016introduction} for
the recursion and \textcite[Sec.~5.4, pp.~136--137]{ljungqvist2018recursive} for
stabilisability and its role in the dual linear-regulator problem, whose Riccati
equation is the same object.

\section{Probability Inequalities}\label{sec:prob-inequalities}

\begin{theorem}[Markov and Chebyshev]\label{thm:markov}
	Let $X\geq0$ be a random variable and $a>0$. Then
	\begin{equation}\label{eq:markov}
		\Prob(X\geq a)\leq\frac{\E[X]}{a} .
	\end{equation}
	Consequently, for any square-integrable $Z$ and $k>0$,
	\begin{equation}\label{eq:chebyshev}
		\Prob\bigl(\abs{Z-\E[Z]}\geq k\bigr)\leq\frac{\Var(Z)}{k^{2}} .
	\end{equation}
\end{theorem}

\begin{proof}
	Since $X\geq0$, the pointwise inequality $X\geq a\one_{[X\geq a]}$ holds:
	on $[X\geq a]$ the left side is at least $a$, and elsewhere the right side is
	zero. Taking expectations and using monotonicity and linearity
	(\autoref{prop:expectation-properties}) gives
	$\E[X]\geq a\,\E[\one_{[X\geq a]}]=a\Prob(X\geq a)$, which is
	\eqref{eq:markov}. For \eqref{eq:chebyshev}, apply \eqref{eq:markov} to the
	non-negative variable $X\eqdef(Z-\E[Z])^{2}$ with threshold $k^{2}$ and use
	$\{\abs{Z-\E[Z]}\geq k\}=\{X\geq k^{2}\}$ together with $\E[X]=\Var(Z)$.
\end{proof}

\begin{remark}[Sharpness and use]\label{rem:markov-sharp}
	Markov's inequality is attained: if $X=a$ with probability $\E[X]/a$ and $X=0$
	otherwise, \eqref{eq:markov} holds with equality, so no distribution-free
	improvement is possible. Its role is to convert a statement about a mean into a
	statement about a tail using no distributional assumption, and this is what
	drives the weak law of large numbers: for i.i.d.\ square-integrable
	$(Z_i)$ with mean $\mu$ and variance $\sigma^{2}$, the sample mean
	$\bar Z_n$ has $\E[\bar Z_n]=\mu$ and, by independence,
	$\Var(\bar Z_n)=\sigma^{2}/n$, so \eqref{eq:chebyshev} gives
	$\Prob(\abs{\bar Z_n-\mu}\geq k)\leq\sigma^{2}/(nk^{2})\to0$ for every $k>0$.
	Convergence in probability of sample averages is therefore a two-line
	consequence of the variance identity and Markov's inequality, requiring no
	central limit theorem and no normality.
\end{remark}

\section{Convergence of Random Variables}\label{sec:prob-convergence}

Everything asymptotic in econometrics --- consistency, asymptotic normality, the
behaviour of a sample average as the sample grows --- is a statement that a
sequence of random variables converges. Unlike convergence in $\R^n$, there are
several inequivalent senses in which it can do so, and the theorems that license
exchanging a limit with an expectation carry hypotheses that are routinely
needed and routinely omitted.

\begin{definition}[Modes of convergence]\label{def:modes-convergence}
	Let $X$ and $X_1,X_2,\dots$ be random variables on $(\Omega,\mathcal{F},\Prob)$.
	\begin{enumerate}[(i)]
		\item $X_n\to X$ \dfnas{almost sure convergence}{almost surely}, written
		      $X_n\overset{\text{a.s.}}{\to}X$, if
		      $\Prob\bigl(\{\omega : X_n(\omega)\to X(\omega)\}\bigr)=1$.
		\item $X_n\to X$ \dfnas{convergence in probability}{in probability}, written
		      $X_n\overset{\Prob}{\to}X$, if $\Prob(\abs{X_n-X}>\varepsilon)\to0$ for
		      every $\varepsilon>0$.
		\item $X_n\to X$ \dfnas{convergence in Lp@convergence in $\Lspace{p}$}{in
			      $\Lspace{p}$}, $p\geq1$, if $\E\bigl[\abs{X_n-X}^{p}\bigr]\to0$.
		      \nidx{Lp}{$\Lspace{p}$, space of $p$-integrable random variables}
		\item $X_n\to X$ \dfnas{convergence in distribution}{in distribution},
		      written $X_n\overset{d}{\to}X$, if $\E[g(X_n)]\to\E[g(X)]$ for every
		      bounded continuous $g\colon\R\to\R$; equivalently, if
		      $F_n(x)\to F(x)$ at every continuity point $x$ of the distribution
		      function $F$ of $X$.
	\end{enumerate}
	\nidx{convasconv}{$\overset{\text{a.s.}}{\to}$, $\overset{\Prob}{\to}$,
		$\overset{d}{\to}$, modes of convergence}
\end{definition}

Only (iv) is a statement about distributions rather than about the random
variables themselves; $X_n$ and $X$ need not even live on the same probability
space for it to make sense. This is why it is the mode in which limit theorems
are stated, and the mode that supports no arithmetic. Let $Z\sim N(0,1)$ and set
$X_n\eqdef Z$, $Y_n\eqdef-Z$ for every $n$, and let $X,Y$ be \emph{independent}
standard normals. Then $X_n\overset{d}{\to}X$ and $Y_n\overset{d}{\to}Y$, since
each side has the same marginal law at every $n$, while $X_n+Y_n=0$ has the
degenerate limit and $X+Y\sim N(0,2)$. The marginal laws leave the joint law
undetermined, and sums are a function of the joint law.

\begin{proposition}[The hierarchy]\label{prop:convergence-hierarchy}
	Let $X,X_1,X_2,\dots$ be random variables and $p\geq1$.
	\begin{enumerate}[(i)]
		\item $\Lspace{p}$ convergence implies convergence in probability.
		\item Almost sure convergence implies convergence in probability.
		\item Convergence in probability implies convergence in distribution.
		\item Convergence in distribution to a constant $c$ implies convergence in
		      probability to $c$.
	\end{enumerate}
	None of the converses of (i)--(iii) holds, and (i) and (ii) imply neither
	the other.
\end{proposition}

\begin{proof}
	\emph{$\Lspace{p}\Rightarrow\Prob$.} By \autoref{thm:markov} applied to
	$\abs{X_n-X}^{p}$,
	$\Prob(\abs{X_n-X}>\varepsilon)\leq\varepsilon^{-p}\E[\abs{X_n-X}^{p}]\to0$.

	\emph{a.s.\ $\Rightarrow\Prob$.} Fix $\varepsilon>0$ and set
	$A_n\eqdef\bigcup_{m\geq n}\{\abs{X_m-X}>\varepsilon\}$. The sets $A_n$
	decrease, and on the event $\{X_n\to X\}$ no $\omega$ lies in
	$\bigcap_nA_n$; hence $\Prob(\bigcap_nA_n)=0$ and, by continuity from above
	(\autoref{prop:probability-basics}(vi)), $\Prob(A_n)\to0$. Since
	$\{\abs{X_n-X}>\varepsilon\}\subseteq A_n$, the claim follows.

	\emph{$\Prob\Rightarrow d$.} Let $g$ be bounded by $M$ and continuous, and fix
	$\varepsilon>0$. Choose $K$ with $\Prob(\abs{X}>K)<\varepsilon$; $g$ is
	uniformly continuous on $[-K-1,K+1]$ by \autoref{thm:heine-cantor}, so there is
	$\delta\in(0,1)$ with $\abs{g(x)-g(y)}<\varepsilon$ whenever
	$\abs{x-y}<\delta$ and $\abs{x}\leq K$. Splitting according to
	$\{\abs{X_n-X}<\delta\}$, $\{\abs{X}>K\}$ and the remainder,
	\[
		\bigl\lvert\E[g(X_n)]-\E[g(X)]\bigr\rvert
		\leq\varepsilon+2M\Prob(\abs{X}>K)+2M\Prob(\abs{X_n-X}\geq\delta)
		\leq3\varepsilon+2M\Prob(\abs{X_n-X}\geq\delta),
	\]
	and the last term vanishes in the limit.

	\emph{Failures of the converses.} On $([0,1],\mathcal{B},\text{Leb})$ let
	$X_n=\one_{[j2^{-k},(j+1)2^{-k}]}$ where $n=2^{k}+j$, $0\leq j<2^{k}$: the
	intervals sweep $[0,1]$ repeatedly with lengths $2^{-k}\to0$, so
	$X_n\overset{\Prob}{\to}0$, while $\limsup_nX_n(\omega)=1$ for every $\omega$,
	so $X_n$ converges almost surely to nothing. Taking instead
	$X_n=n\one_{[0,1/n]}$ gives $X_n\overset{\text{a.s.}}{\to}0$ with
	$\E[X_n]=1$ for every $n$, so there is no $\Lspace{1}$ convergence. And if
	$X\sim N(0,1)$ and $X_n\eqdef-X$ for every $n$, then
	$X_n\overset{d}{\to}X$ by symmetry while $\abs{X_n-X}=2\abs{X}$ does not
	converge to $0$ in probability.

	\emph{Convergence in distribution to a constant.} If $X_n\overset{d}{\to}c$
	then for $\varepsilon>0$ the distribution function of $X_n$ tends to $0$ below
	$c-\varepsilon$ and to $1$ above $c+\varepsilon$, both continuity points of the
	degenerate limit, so $\Prob(\abs{X_n-c}>\varepsilon)\to0$.
\end{proof}

\begin{remark}[Where each mode is used]\label{rem:modes-econometrics}
	Consistency of an estimator is convergence in probability of
	$\hat\theta_n$ to $\theta_0$; the notation $\plim\hat\theta_n=\theta_0$ is
	\nidx{plim}{$\plim$, probability limit}
	\autoref{def:modes-convergence}(ii). Asymptotic normality is convergence in
	distribution of $\sqrt n(\hat\theta_n-\theta_0)$. Mean-square consistency is
	convergence in $\Lspace{2}$, and it is what the projection argument of
	\autoref{ch:projection} delivers directly, since
	$\E[\abs{X_n-X}^{2}]$ is a squared distance in the Hilbert space
	$\Lspace{2}$. The hierarchy of \autoref{prop:convergence-hierarchy} orders
	the first two: mean-square consistency implies consistency and the converse
	fails. Asymptotic normality does not belong to that chain, being a statement
	about a different sequence. Consistency concerns $\hat\theta_n-\theta_0$;
	asymptotic normality concerns the rescaled error
	$\sqrt n(\hat\theta_n-\theta_0)$. If the rescaled error converges in
	distribution at all, then
	$\hat\theta_n-\theta_0=n^{-1/2}\bigl[\sqrt n(\hat\theta_n-\theta_0)\bigr]
		\overset{d}{\to}0$ by \autoref{thm:slutsky}(ii), hence
	$\overset{\Prob}{\to}0$ by \autoref{prop:convergence-hierarchy}(iv). So
	asymptotic normality implies consistency, and supplies in addition the rate
	$n^{-1/2}$ and the limiting distribution. What mean-square consistency
	supplies beyond either is convergence of second moments, which a
	distributional limit alone does not give: $X_n=n\one_{[0,1/n]}$ converges to
	$0$ in distribution and in probability with $\E[X_n]=1$ throughout.
\end{remark}

\begin{theorem}[Continuous mapping, Slutsky, delta
		method]\label{thm:slutsky}
	Let $X_n,X$ take values in $\R^{k}$ and let $c\in\R^{k}$, $a\in\R$.
	\begin{enumerate}[(i)]
		\item \emph{(Continuous mapping.)} If $X_n\overset{d}{\to}X$ and
		      $h\colon\R^{k}\to\R^{m}$ is continuous on a set $C$ with
		      $\Prob(X\in C)=1$, then $h(X_n)\overset{d}{\to}h(X)$. The same
		      implication holds with $\overset{d}{\to}$ replaced throughout by
		      $\overset{\Prob}{\to}$ or by almost sure convergence.
		\item \emph{(Slutsky.)} Let $X_n\overset{d}{\to}X$ in $\R^{k}$,
		      $Y_n\overset{\Prob}{\to}c$ in $\R^{k}$, and $A_n\overset{\Prob}{\to}A$
		      for random $m\times k$ matrices $A_n$ and a constant matrix $A$. Then
		      $X_n+Y_n\overset{d}{\to}X+c$ and $A_nX_n\overset{d}{\to}AX$. If
		      $k=1$ and $a\neq0$ with $Y_n\overset{\Prob}{\to}a$ real, then
		      $X_n/Y_n\overset{d}{\to}X/a$.
		\item \emph{(Delta method.)} If
		      $\sqrt n(X_n-a)\overset{d}{\to}X$ with $X_n,a\in\R^{k}$ and
		      $h\colon\R^{k}\to\R^{m}$ is differentiable at $a$, then
		      $\sqrt n\bigl(h(X_n)-h(a)\bigr)\overset{d}{\to}Dh(a)X$.
	\end{enumerate}
\end{theorem}

\begin{proof}
	Throughout, convergence in distribution is used in the form
	$\E[g(X_n)]\to\E[g(X)]$ for every bounded continuous $g$, as in
	\autoref{def:modes-convergence}(iv).

	\emph{A tightness bound.} If $X_n\overset{d}{\to}X$ then for every
	$\varepsilon>0$ there is $K$ with
	$\limsup_n\Prob(\norm{X_n}>K+1)\leq\varepsilon$. Indeed
	$h_K(x)\eqdef\min\{1,\max\{0,\norm{x}-K\}\}$ is bounded and continuous with
	$\one_{\{\norm{x}>K+1\}}\leq h_K(x)\leq\one_{\{\norm{x}>K\}}$, so
	$\limsup_n\Prob(\norm{X_n}>K+1)\leq\lim_n\E[h_K(X_n)]=\E[h_K(X)]
		\leq\Prob(\norm{X}>K)$, and the right side tends to $0$ as $K\to\infty$.

	\emph{(ii), sums.} Let $g$ be continuous with $\abs{g}\leq M$ and fix
	$\varepsilon>0$. Choose $K$ as above. The closed ball of radius
	$K+1+\norm{c}+1$ is compact, so $g$ is uniformly continuous on it by
	\autoref{thm:heine-cantor}: there is $\delta\in(0,1)$ with
	$\abs{g(x)-g(y)}<\varepsilon$ for $x,y$ in that ball with
	$\norm{x-y}<\delta$. On the event
	$\{\norm{Y_n-c}<\delta\}\cap\{\norm{X_n}\leq K+1\}$ the two points
	$X_n+Y_n$ and $X_n+c$ lie in the ball and differ by less than $\delta$, so
	\[
		\bigl\lvert\E[g(X_n+Y_n)]-\E[g(X_n+c)]\bigr\rvert
		\leq\varepsilon+2M\bigl[\Prob(\norm{Y_n-c}\geq\delta)
			+\Prob(\norm{X_n}>K+1)\bigr].
	\]
	The first probability vanishes and the second has $\limsup\leq\varepsilon$.
	Since $x\mapsto g(x+c)$ is bounded and continuous,
	$\E[g(X_n+c)]\to\E[g(X+c)]$, and therefore
	$\limsup_n\abs{\E[g(X_n+Y_n)]-\E[g(X+c)]}\leq\varepsilon(1+2M)$ for every
	$\varepsilon>0$.

	\emph{(ii), products and quotients.} For the matrix product the same
	argument applies on the event
	$\{\norm{A_n-A}_{2}<\delta/(K+1)\}\cap\{\norm{X_n}\leq K+1\}$, on
	which $\norm{A_nX_n-AX_n}\leq\norm{A_n-A}_{2}\norm{X_n}<\delta$,
	with the ball of radius $(\norm{A}_{2}+1)(K+1)$ and with
	$\E[g(AX_n)]\to\E[g(AX)]$ by (i). If $a\neq0$ then
	$1/Y_n\overset{\Prob}{\to}1/a$, since $t\mapsto1/t$ is continuous at $a$ and
	convergence in probability is preserved under a map continuous at the limit
	point; the product case then gives $X_n/Y_n\overset{d}{\to}X/a$. Degeneracy
	of the second limit is what makes these arguments work, and the failure
	recorded before \autoref{prop:convergence-hierarchy} shows that it cannot be
	dropped.

	\emph{(iii).} Differentiability at $a$ gives
	$h(x)-h(a)=Dh(a)(x-a)+\norm{x-a}\psi(x)$ with $\psi(a)\eqdef0$ and
	$\psi(x)\to0$ as $x\to a$. Applying (ii) with the constant factor
	$n^{-1/2}\to0$ gives $X_n-a\overset{d}{\to}0$, hence
	$X_n-a\overset{\Prob}{\to}0$ by
	\autoref{prop:convergence-hierarchy}(iv). For $\varepsilon>0$ pick $\eta>0$
	with $\norm{\psi(x)}<\varepsilon$ whenever $\norm{x-a}<\eta$; then
	$\Prob(\norm{\psi(X_n)}\geq\varepsilon)\leq\Prob(\norm{X_n-a}\geq\eta)\to0$,
	so $\psi(X_n)\overset{\Prob}{\to}0$. The tightness bound applied to
	$\sqrt n(X_n-a)$ makes $\norm{\sqrt n(X_n-a)}$ bounded in probability, so
	the product $\sqrt n\norm{X_n-a}\psi(X_n)\overset{\Prob}{\to}0$. Finally
	$\sqrt n\bigl(h(X_n)-h(a)\bigr)
		=Dh(a)\sqrt n(X_n-a)+\sqrt n\norm{X_n-a}\psi(X_n)$, where the first term
	converges in distribution to $Dh(a)X$ by (i) applied to the continuous map
	$Dh(a)$, and (ii) absorbs the second.

	Part (i) for convergence in distribution is
	\textcite[Thm.~9.4.3, p.~562]{corbae2009introduction}; its proof runs
	through the portmanteau characterisation of weak convergence, which is not
	developed here. For convergence in probability and almost sure convergence
	the statement is elementary: almost sure convergence passes through $h$
	pointwise on the set where $X\in C$, and the probability version follows by
	the subsequence criterion.
\end{proof}

\subsection{Interchanging limits and expectations}\label{subsec:prob-convergence-thms}

\begin{theorem}[Monotone convergence, Fatou, dominated
		convergence]\label{thm:convergence-theorems}
	Let $X,X_1,X_2,\dots$ be random variables on $(\Omega,\mathcal{F},\Prob)$.
	\begin{enumerate}[(i)]
		\item \emph{(Monotone convergence.)} If $0\leq X_n\leq X_{n+1}$ for every $n$
		      and $X_n\to X$ pointwise, then $\E[X_n]\to\E[X]$, both sides being
		      allowed the value $+\infty$.
		\item \emph{(Fatou.)} If $X_n\geq0$ for every $n$, then
		      $\E[\liminf_nX_n]\leq\liminf_n\E[X_n]$.
		\item \emph{(Dominated convergence.)} If $X_n\to X$ almost surely and there is
		      an integrable $Y$ with $\abs{X_n}\leq Y$ almost surely for every $n$,
		      then $X$ and every $X_n$ are integrable and $\E[X_n]\to\E[X]$; moreover
		      $\E[\abs{X_n-X}]\to0$.
	\end{enumerate}
\end{theorem}

These are Theorems~2.4 (p.~21), 2.8 (p.~25) and 2.11 (p.~30) of
\textcite{legall2022measure}, stated there for a general measure space; the proofs
are omitted here because they belong to the construction of the Lebesgue
integral, which \autoref{ch:integration} deliberately does not carry out, and
because nothing in this volume depends on their internal structure. Each is used
below only as a licence to exchange a limit with an integral.

\begin{remark}[Where each hypothesis is used]\label{rem:convergence-hypotheses}
	Each of (i)--(iii) fails without its hypothesis, and the same example refutes
	all three. On $([0,1],\mathcal{B},\text{Leb})$ put $X_n=n\one_{[0,1/n]}$. Then
	$X_n\to0$ pointwise, so $\E[\lim X_n]=0$, while $\E[X_n]=1$ for every $n$:
	monotonicity fails in (i), Fatou's inequality in (ii) is strict, and no
	integrable dominating $Y$ exists in (iii), since $\sup_nX_n\geq1/x$ on
	$(0,1]$, which is not integrable. In economic terms the example is a sequence
	of lotteries whose payoffs grow as fast as their probabilities shrink; expected
	utility is continuous along such sequences only under a domination or
	uniform-integrability condition, and this is what a bounded utility function
	buys.

	Part (iii) is what lets expectations pass through limits in
	\autoref{thm:cond-exp-exists} and in \autoref{prop:cond-exp-density}, and
	\autoref{thm:differentiation-under-integral} is the form of it needed to
	differentiate under an integral sign --- the step behind every comparative
	static of an expected-utility problem.
\end{remark}

\subsection{Limit theorems}\label{subsec:prob-limit-theorems}

\begin{theorem}[Laws of large numbers]\label{thm:lln}
	Let $(X_n)_{n\in\N}$ be independent and identically distributed with
	$\mu\eqdef\E[X_1]$, and let $\bar X_n\eqdef n^{-1}\sum_{i=1}^{n}X_i$.
	\begin{enumerate}[(i)]
		\item \emph{(Weak law.)} If $\Var(X_1)=\sigma^{2}<\infty$, then
		      $\E[(\bar X_n-\mu)^{2}]=\sigma^{2}/n\to0$, so
		      $\bar X_n\to\mu$ in $\Lspace{2}$ and hence in probability.
		\item \emph{(Strong law.)} If $\E\abs{X_1}<\infty$, then
		      $\bar X_n\overset{\text{a.s.}}{\to}\mu$.
	\end{enumerate}
\end{theorem}

\begin{proof}
	(i) $\E[\bar X_n]=\mu$ by linearity, and independence gives
	$\Var(\bar X_n)=n^{-2}\sum_i\Var(X_i)=\sigma^{2}/n$ by
	\autoref{prop:variance-independent}. Convergence in probability follows from
	\autoref{prop:convergence-hierarchy}, or directly from \eqref{eq:chebyshev} as
	in \autoref{rem:markov-sharp}.

	(ii) is Theorem~10.8 (p.~206) of \textcite{legall2022measure}. The proof rests
	on Kolmogorov's zero-one law and a truncation argument, neither of which is
	developed here; what matters below is only the conclusion, and the hypothesis
	$\E\abs{X_1}<\infty$ is not merely sufficient but necessary for the limit to be
	defined.
\end{proof}

\begin{theorem}[Central limit theorem]\label{thm:clt}
	Let $(X_n)_{n\in\N}$ be independent and identically distributed with
	$\E[X_1]=\mu$ and $\Var(X_1)=\sigma^{2}\in(0,\infty)$. Then
	\begin{equation}\label{eq:clt}
		\sqrt n\,\bigl(\bar X_n-\mu\bigr)\overset{d}{\longrightarrow}N(0,\sigma^{2}).
	\end{equation}
\end{theorem}

This is Theorem~10.15 (p.~219) of \textcite{legall2022measure}; the proof uses
characteristic functions and L\'evy's continuity theorem, neither of which is
introduced in this volume, and is not reproduced. Two features of the statement
are what applied work relies on. The scaling is $\sqrt n$ and not $n$: by
\autoref{thm:lln}(i) the deviation $\bar X_n-\mu$ has standard deviation
$\sigma/\sqrt n$, so $\sqrt n$ is the only normalisation under which the limit is
non-degenerate. And the limit law does not depend on the distribution of $X_1$
beyond its first two moments, which is why a standard error and a normal quantile
suffice to form a confidence interval without knowing the data-generating
distribution.

\begin{remark}[Asymptotic normality of least squares]\label{rem:ols-asymptotics}
	Return to the estimator $\hat\beta=(X^{\transp}X)^{-1}X^{\transp}y$ of
	\autoref{prop:ols-projection} with $y=X\beta+\varepsilon$. Then
	\[
		\sqrt n\,(\hat\beta-\beta)
		=\Bigl(\tfrac1nX^{\transp}X\Bigr)^{-1}\cdot
		\tfrac{1}{\sqrt n}X^{\transp}\varepsilon .
	\]
	The hypotheses are that the pairs $(x_i,\varepsilon_i)$ are i.i.d.\ and
	\begin{equation}\label{eq:ols-asy-assumptions}
		\E[x_i\varepsilon_i]=0,\qquad
		Q\eqdef\E[x_ix_i^{\transp}]\ \text{finite and positive definite},\qquad
		\Omega\eqdef\E\bigl[\varepsilon_i^{2}x_ix_i^{\transp}\bigr]\ \text{finite},
	\end{equation}
	the last of which is the statement $\E[\varepsilon_i^{2}\norm{x_i}^{2}]<\infty$
	up to a change of norm. Finiteness of $Q$ makes the entries of
	$x_ix_i^{\transp}$ integrable, so \autoref{thm:lln} gives
	$n^{-1}X^{\transp}X\to Q$ almost surely; matrix inversion is continuous at
	every invertible matrix, so positive definiteness of $Q$ and
	\autoref{thm:slutsky}(i) give
	$(n^{-1}X^{\transp}X)^{-1}\to Q^{-1}$ almost surely, hence in probability.
	Finiteness of $\Omega$ makes the i.i.d.\ vectors $x_i\varepsilon_i$
	square-integrable with mean zero, so \autoref{thm:clt} gives
	$n^{-1/2}X^{\transp}\varepsilon\overset{d}{\to}N(0,\Omega)$. Now
	\autoref{thm:slutsky}(ii), in the form that a random matrix converging in
	probability to a constant matrix may be substituted into a product, yields
	$\sqrt n(\hat\beta-\beta)\overset{d}{\to}Q^{-1}N(0,\Omega)
		=N(0,Q^{-1}\Omega Q^{-1})$, the sandwich formula. Dropping positive
	definiteness of $Q$ breaks the argument at the inversion step, and dropping
	finiteness of $\Omega$ breaks it at the central limit theorem, leaving
	$\hat\beta$ consistent but with no $\sqrt n$ limit. Invertibility of $Q$ is
	the population version of the
	full-rank condition, and \autoref{rem:collinearity} says that when $Q$ is nearly
	singular the asymptotic variance along the least-excited direction of the
	regressors is large: $\kappa(Q)$ measures how much information the design
	carries about the worst-identified linear combination of coefficients.
\end{remark}

\subsection{Martingales}\label{subsec:prob-martingales}

\begin{definition}[Filtration and martingale]\label{def:martingale}
	A \dfnas{filtration}{filtration} on $(\Omega,\mathcal{F},\Prob)$ is an
	increasing sequence $\mathcal{F}_0\subseteq\mathcal{F}_1\subseteq\dots
		\subseteq\mathcal{F}$ of $\sigma$-algebras. A process $(M_n)_{n\geq0}$ is
	\dfnas{adapted process}{adapted} if each $M_n$ is $\mathcal{F}_n$-measurable,
	and is a \dfnas{martingale}{martingale} if it is adapted, $\E\abs{M_n}<\infty$
	for every $n$, and
	\begin{equation}\label{eq:martingale}
		\E\bigl[M_{n+1}\mid\mathcal{F}_n\bigr]=M_n\qquad\text{for every }n\geq0 .
	\end{equation}
	It is a \dfnas{supermartingale}{supermartingale} if $\leq$ holds in
	\eqref{eq:martingale} and a \dfnas{submartingale}{submartingale} if $\geq$
	does. A \dfnas{stopping time}{stopping time} is a random
	$T\colon\Omega\to\N_0\cup\{+\infty\}$ with $\{T\leq n\}\in\mathcal{F}_n$ for
	every $n$.
	\nidx{filtration}{$\mathcal{F}_n$, filtration}
\end{definition}

The filtration is the formalisation of information arriving over time, and
$\mathcal{F}_n$ is the $\sigma$-algebra of \autoref{def:conditional-expectation-sigma}
representing what is known at date $n$. A stopping time is a rule for quitting
that uses only information available when it is exercised: the condition
$\{T\leq n\}\in\mathcal{F}_n$ excludes rules such as ``sell at the peak''.

\begin{proposition}[Two constructions]\label{prop:martingale-examples}
	\begin{enumerate}[(i)]
		\item \emph{(Doob martingale.)} If $\E\abs{Z}<\infty$ then
		      $M_n\eqdef\E[Z\mid\mathcal{F}_n]$ is a martingale.
		\item \emph{(Random walk.)} If $(\xi_n)_{n\geq1}$ are independent with
		      $\E[\xi_n]=0$ and $\E\abs{\xi_n}<\infty$, and
		      $\mathcal{F}_n=\sigma(\xi_1,\dots,\xi_n)$, then
		      $M_n\eqdef\sum_{i\leq n}\xi_i$ is a martingale.
		\item If $(M_n)$ is a martingale and $\phi$ is convex with
		      $\E\abs{\phi(M_n)}<\infty$, then $(\phi(M_n))$ is a submartingale.
	\end{enumerate}
\end{proposition}

\begin{proof}
	(i) Integrability is \autoref{thm:cond-exp-exists}, and
	$\E[M_{n+1}\mid\mathcal{F}_n]=\E\bigl[\E[Z\mid\mathcal{F}_{n+1}]
		\mid\mathcal{F}_n\bigr]=\E[Z\mid\mathcal{F}_n]=M_n$ by the tower property
	\eqref{eq:lie}, which applies because $\mathcal{F}_n\subseteq\mathcal{F}_{n+1}$.

	(ii) $\E[M_{n+1}\mid\mathcal{F}_n]=M_n+\E[\xi_{n+1}\mid\mathcal{F}_n]
		=M_n+\E[\xi_{n+1}]=M_n$, using independence of $\xi_{n+1}$ from
	$\mathcal{F}_n$.

	(iii) By conditional Jensen (\autoref{prop:conditional-jensen}),
	$\E[\phi(M_{n+1})\mid\mathcal{F}_n]\geq\phi\bigl(\E[M_{n+1}\mid\mathcal{F}_n]\bigr)
		=\phi(M_n)$.
\end{proof}

\begin{theorem}[Optional stopping]\label{thm:optional-stopping}
	Let $(M_n)$ be a martingale and $T$ a stopping time. Then the stopped process
	$(M_{n\wedge T})_{n\geq0}$ is a martingale. If in addition $T$ is bounded, then
	$M_T$ is integrable and $\E[M_T]=\E[M_0]$.
\end{theorem}

This is Theorem~12.12 (p.~265) of \textcite{legall2022measure}, where it is
obtained by writing $M_{n\wedge T}=M_0+(H\bullet M)_n$ for the predictable
sequence $H_n=\one_{\{T\geq n\}}$ and using that a predictable transform of a
martingale is a martingale; the argument is short and is not reproduced. The
boundedness of $T$ cannot be dropped: for the symmetric random walk with
$M_0=0$ and $T\eqdef\inf\{n : M_n=1\}$ one has $T<\infty$ almost surely and
$\E[M_T]=1\neq0=\E[M_0]$. The gambler who plays until ahead is certain to
succeed and has unbounded expected playing time; the failure of the conclusion at
an unbounded stopping time is what rules out that strategy in a model with
borrowing limits.

\begin{remark}[Martingales in the models of this volume]\label{rem:martingale-econ}
	Three results already proved acquire their standard names once
	\eqref{eq:martingale} is available.

	The tower property \eqref{eq:lie} says that the conditional expectation of a
	fixed random variable given a growing information set is a martingale, by
	\autoref{prop:martingale-examples}(i). Applied to a forecast, this is the
	statement that revisions to a rational forecast are unpredictable from current
	information --- the testable content of rational expectations, and the reason a
	regression of the forecast revision on any date-$n$ variable should have zero
	coefficient.

	In \autoref{thm:ftap}, absence of arbitrage in the one-period model produced a
	strictly positive state price vector $\psi$. Write
	$R_0\eqdef1/\sum_s\psi_s$ for the gross return on the riskless claim and
	$\Prob^{\ast}(s)\eqdef\psi_sR_0$, which is a probability measure because its
	weights are positive and sum to one. Then \eqref{eq:state-prices} becomes
	$q_j=\E^{\ast}[R_{\cdot j}]/R_0$: every asset's price is its expected payoff
	under $\Prob^{\ast}$, discounted at the riskless rate. Applying this at each
	date of a multi-period model, with $\mathcal{F}_n$ the information available at
	date $n$, gives $S_n=\E^{\ast}[S_{n+1}\mid\mathcal{F}_n]/R_0$, so the discounted
	price process $R_0^{-n}S_n$ is a martingale under $\Prob^{\ast}$. The
	multi-period statement is not proved here --- it requires the one-period result
	at every node, together with a consistency condition tying the measures across
	dates --- but the algebra at a single date is \eqref{eq:state-prices} rewritten,
	and $\Prob^{\ast}$ is the risk-neutral measure.

	In the Lucas model of \autoref{sec:dyn-lucas}, the Euler equation
	$p_t u'(c_t)=\beta\E_t\bigl[(p_{t+1}+d_{t+1})u'(c_{t+1})\bigr]$ says that
	$\beta^{t}u'(c_t)p_t$ plus accumulated discounted dividends is a martingale;
	the transversality condition is the requirement that this martingale not
	diverge, which is how \autoref{rem:transversality} rules out bubbles.
\end{remark}

\section{Chapter Summary}\label{sec:prob-summary}

A probability space consists of a sample space, a $\sigma$-algebra of events
closed under complementation and countable intersection, and a countably additive
measure of total mass one; monotonicity, finite additivity and
inclusion--exclusion follow from the two axioms. A random variable is a
measurable real function on the sample space, its expectation the mass-weighted
sum of its values, and the variance identity
$\Var(X)=\E[X^{2}]-(\E[X])^{2}$ is the special case, at $c=0$, of the
decomposition $\E[(X-c)^{2}]=\Var(X)+(\E[X]-c)^{2}$ that identifies the mean as
the best constant predictor. Jensen's inequality follows from the existence of a
supporting line to a convex function, and yields the sign of the risk premium
without differentiability.

Conditional probability makes $\Prob(\cdot\mid H)$ a probability measure in its
own right; Bayes' rule reverses the direction of conditioning, and in odds form
states that posterior odds are prior odds times the likelihood ratio. The
denominator, computed by the law of total probability, is what the base-rate
fallacy neglects. Conditional expectation given a random variable is itself a
random variable, satisfies the pull-out property
$\E[g(X)Y\mid X]=g(X)\E[Y\mid X]$ and the law of iterated expectations
$\E[\E[X\mid Y]]=\E[X]$, and, viewed in $L^{2}$, is the orthogonal projection
of $Y$ onto the functions of $X$. That identification supplies at once the
best-predictor property, the law of total variance, and the fact that the best
linear predictor is a strictly worse projection onto a strictly smaller subspace,
the two coinciding exactly in the jointly normal case.

Gaussian updating is conjugate: posterior precision is the sum of prior and
signal precisions, independent of the realisation, and the posterior mean is the
corresponding precision-weighted average, linear in the signal. With many
independent signals the same arithmetic applies to their sum, and the sequential
and simultaneous updates agree. Applied to a CARA investor the linearity yields
demand $(\mu-p)/(a\sigma^{2})$ with no wealth effect, hence prices that are linear
statistics of dispersed information; applied recursively it yields the Kalman
recursion, whose variance obeys the Riccati equation shown in
\autoref{ch:completeness} to be a contraction. Markov's inequality bounds a tail
by a mean with no distributional assumption, and through Chebyshev delivers the
weak law of large numbers.

\begin{exercise}\label{ex:prob-1}
	Prove the general Bayes formula \eqref{eq:bayes-partition} directly from
	\autoref{def:conditional-probability} without invoking
	\eqref{eq:bayes-1}, and verify it against \autoref{eg:poll}. Then show that the
	odds form \eqref{eq:odds-form} extends to a partition of size $n$ as the
	statement that $\Prob(H_i\mid D)/\Prob(H_j\mid D)$ equals the prior odds
	ratio times the likelihood ratio, so that the normalising constant never has to
	be computed to compare two hypotheses.
\end{exercise}

\begin{exercise}\label{ex:prob-2}
	Prove the law of total variance \eqref{eq:variance-decomposition} from
	\autoref{thm:best-predictor} and \autoref{thm:lie}. Then show that
	$\Var\bigl(\E[Y\mid X]\bigr)\leq\Var(Y)$, with equality if and only if $Y$ is
	almost surely a function of $X$, and interpret the ratio of the two as the
	population $R^{2}$ of the regression of $Y$ on $X$.
\end{exercise}

\begin{exercise}\label{ex:prob-3}
	In the setting of \autoref{prop:microstructure}, compute
	$\Prob(V=H\mid B)$ after a buy, and show that the market maker's bid
	$\E[V\mid S]$ and ask $\E[V\mid B]$ satisfy
	$\E[V\mid S]<\E[V]<\E[V\mid B]$ with the spread strictly increasing in $q$
	and vanishing at $q=0$. Then verify that the martingale property
	$\E\bigl[\E[V\mid\text{trade}]\bigr]=\E[V]$ holds, and explain why this is
	\autoref{thm:lie} rather than an assumption about the market maker's behaviour.
\end{exercise}
